% concatenated from main.tex, preamble.tex
\documentclass{article}


% if you need to pass options to natbib, use, e.g.:
    \PassOptionsToPackage{numbers, sort&compress}{natbib}

% ready for submission
\usepackage[final]{neurips_2025}

\usepackage[utf8]{inputenc} % allow utf-8 input
\usepackage[T1]{fontenc}    % use 8-bit T1 fonts
\usepackage{hyperref}       % hyperlinks
\usepackage{url}            % simple URL typesetting
\usepackage{booktabs}       % professional-quality tables
\usepackage{amsfonts}       % blackboard math symbols
\usepackage{nicefrac}       % compact symbols for 1/2, etc.
\usepackage{microtype}      % microtypography
% \usepackage{xcolor}         % colors

\usepackage{enumitem}

%Packages 
\usepackage{amsmath, amsthm, amsfonts, amssymb}
\usepackage{thmtools}
\usepackage{mathtools}
\usepackage{graphicx}
\graphicspath{{./Figures/}}
\usepackage{subfigure}
\usepackage{algorithm}
\usepackage{algpseudocode}
\usepackage{booktabs, multirow, array}
\usepackage[table]{xcolor}
%Define colors
\definecolor{darkblue}{RGB}{0,0,180}  % a very dark blue


% \hypersetup{hidelinks}
\hypersetup{
  colorlinks=true,     % use colored text instead of boxes
  linkcolor=darkblue,     % color for internal links (sections, equations)
  citecolor=blue,     % color for bibliographic citations
  urlcolor=blue        % color for external URLs
}

%Macro
\def\cA{\mathcal{A}} \def\cB{\mathcal{B}} \def\cC{\mathcal{C}} \def\cD{\mathcal{D}}
\def\cE{\mathcal{E}} \def\cF{\mathcal{F}} \def\cG{\mathcal{G}} \def\cH{\mathcal{H}}
\def\cI{\mathcal{I}} \def\cJ{\mathcal{J}} \def\cK{\mathcal{K}} \def\cL{\mathcal{L}}
\def\cM{\mathcal{M}} \def\cN{\mathcal{N}} \def\cO{\mathcal{O}} \def\cP{\mathcal{P}}
\def\cQ{\mathcal{Q}} \def\cR{\mathcal{R}} \def\cS{\mathcal{S}} \def\cT{\mathcal{T}}
\def\cU{\mathcal{U}} \def\cV{\mathcal{V}} \def\cW{\mathcal{W}} \def\cX{\mathcal{X}}
\def\cY{\mathcal{Y}} \def\cZ{\mathcal{Z}}


\def\R{\mathbb{R}}
\def\N{\mathbb{N}}
\def\C{\mathbb{C}}
\def\E{\mathbb{E}}
\def\st{\text{s.t.}}
%for feasibility seeking paper
\def\FS{\mathsf{FS}}
\def\tFS{\widetilde{\mathsf{FS}}}
\def\tL{\Tilde{L}}
\def\tgrad{\Tilde{\nabla}}


\newcommand{\bigO}{\mathcal{O}}
\newcommand{\e}[1]{e{#1}}
%For table 
\newcolumntype{C}[1]{>{\centering\arraybackslash}m{#1}}
\newcommand{\thinrule}{\midrule[0.1pt]}

\newcommand{\std}[1]{{\tiny$\pm${#1}}}


%Theorems, corollaries, etc
\newtheorem{theorem}{Theorem}
\newtheorem{corollary}{Corollary}[theorem]
\newtheorem{lemma}{Lemma}
\newtheorem{prop}{Proposition}
\newtheorem{remark}{Remark}
\newtheorem{definition}{Definition}
\newtheorem{assumption}{Assumption}
\usepackage{thmtools} %for restatable theorems

%Operators
% \DeclareMathOperator{\diag}{diag}
% \DeclareMathOperator{\relu}{ReLU}
% \DeclareMathOperator*{\E}{\mathbb{E}}
% \DeclareMathOperator*{\argmax}{arg\,max}
% \DeclareMathOperator*{\argmin}{arg\,min}

%%% For variables 
\def\x{x}
\def\y{y_{\theta}}   \def\hy{\hat{y}_{\theta}}    \def\ty{\Tilde{y}}
\def\z{z}

\def\i{{(i)}}
\def\trun{\mathrm{trun}}


%%% Optimization and linear algebra
\DeclareMathOperator*{\maximize}{maximize}
\DeclareMathOperator{\prox}{prox}
\DeclareMathOperator*{\minimize}{minimize}
\DeclareMathOperator*{\subjectto}{subject\;to}
% \DeclareMathOperator*{\st}{s.t.}
\DeclareMathOperator*{\diag}{diag}
\DeclareMathOperator*{\mean}{mean}
\DeclareMathOperator*{\ReLU}{ReLU}


\newenvironment{ColorBlock}[1]{%
  \rowcolor{#1}%
}{%
  \ignorespacesafterend
}

\newcommand{\todo}[1]{{\color{red} TODO: #1}}
\newcommand{\revise}[1]{{\color{blue} #1}}



\title{FSNet: Feasibility-Seeking Neural Network for Constrained Optimization with Guarantees}

\author{%
  Hoang T. Nguyen \\
  Massachusetts Institute of Technology \\
  Cambridge, MA, USA \\
  \texttt{hoangh@mit.edu} \\
  % examples of more authors
  \And
  Priya L. Donti  \\
  Massachusetts Institute of Technology \\
  Cambridge, MA, USA \\
  \texttt{donti@mit.edu} \\
}


\begin{document}


\maketitle


\begin{abstract}
  Efficiently solving constrained optimization problems is crucial for numerous real-world applications, yet traditional solvers are often computationally prohibitive for real-time use. 
  Machine learning-based approaches have emerged as a promising alternative to provide approximate solutions at faster speeds, but they struggle to strictly enforce constraints, leading to infeasible solutions in practice. 
  To address this, we propose the Feasibility-Seeking Neural Network (FSNet), which integrates a feasibility-seeking step directly into its solution procedure to ensure constraint satisfaction. 
  This feasibility-seeking step solves an unconstrained optimization problem that minimizes constraint violations in a differentiable manner, enabling end-to-end training and providing guarantees on feasibility and convergence. 
  Our experiments across a range of different optimization problems, including both smooth/nonsmooth and convex/nonconvex problems, demonstrate that FSNet can provide feasible solutions with solution quality comparable to (or in some cases better than) traditional solvers, at significantly faster speeds.\footnote{The code is available at \url{https://github.com/MOSSLab-MIT/FSNet}}
  % FSNet thus bridges the gap between performance and reliability, making NN-based optimization practical and trustworthy for real-world, safety-critical applications.
\end{abstract}


\section{Introduction}
Machine learning (ML) has emerged as a promising technique to provide fast surrogate solvers for parametric optimization problems \cite{van2025optimization, amos2023tutorial}.
Specifically, given a set of optimization problems 
that share a fixed structure 
% whose objective and constraints have fixed functional forms
but whose parameters vary between instances, ML can be used to learn approximate, fast-to-evaluate mappings from parameters to optimization solutions.
% \revise{
% This problem aims to minimize an objective function subject to equality and inequality constraints, and the objective and constraint functions depend on parameters but remain the same structure.}
With this approach, many expensive real-world optimization problems can now be solved using ML
in a fraction of a second, with promising examples in electric power systems \cite{park2023compact, khaloie2025review}, building energy management systems \cite{drgovna2021physics}, communication networks \cite{hu2021reconfigurable}, robotic planning and control \cite{wu2024deep, cauligi2020learning}, and supply chain management \cite{akhlaghi2025propel}.
Unfortunately, although fast at producing solutions, ML approaches often struggle to satisfy problem constraints, resulting in solutions that are infeasible and potentially unusable, particularly in safety-critical applications.
While several constraint-enforcement approaches have been proposed to mitigate this problem, including penalty methods \cite{fioretto2021lagrangian}, projection-based methods \cite{min2024hard, huang2021differentiable}, and iterative methods \cite{donti2021dc3, lastrucci2025enforce}, these methods often suffer from poor solution quality, expensive computational costs, and/or lack of guarantees.

To address these challenges, we propose the \textbf{F}easibility-\textbf{S}eeking Neural \textbf{N}etwork (\textbf{FSNet}), a general framework for solving constrained optimization problems using neural networks (NNs). 
Our key idea is to incorporate a differentiable feasibility-seeking step into the NN's training and inference processes to minimize constraint violations while still enabling end-to-end training.
In particular, this feasibility-seeking step uses the NN's prediction as the starting point for a constraint violation minimization problem, which is solved using an iterative optimization solver to ultimately produce
a feasible solution. 
This solution is evaluated via an unsupervised loss function that captures the objective function of the target optimization problem and the distance between the pre- and post-feasibility-seeking points.
The full procedure is trained end-to-end, including backpropagation through the unrolled iterations of the optimization solver used in the feasibility-seeking step.

Our main contributions are summarized as follows: 
\begin{itemize}
    \item \textbf{General framework for constraint enforcement in NNs.} We introduce a principled and general framework, FSNet, that enables NNs to predict solutions satisfying both equality and inequality constraints in continuous optimization problems. The framework accommodates both convex and nonconvex constraints, making it applicable to many problem domains. %applications.

    \item \textbf{Theoretical guarantees on feasibility and convergence.} Under boundedness, $L$-smoothness, and the Polyak–Łojasiewicz (PL) condition, FSNet guarantees exponentially fast convergence to a feasible point and provably reaches a stationary point in the parameter space when being trained using stochastic gradient descent, even when backpropagation is truncated.
    We also specify conditions under which FSNet's predictions provably coincide with a true minimizer of the original constrained optimization problem.
        
    \item \textbf{Strong empirical performance across various classes of problems.} 
    Across convex and nonconvex problems, FSNet produces feasible solutions and achieves
    small optimality gaps (less than 0.2\% in convex cases), while providing solutions orders-of-magnitude faster than traditional solvers.
    % thousands of problem instances 
    % in under half a second. %0.5 seconds. 
    In nonconvex cases, we show that FSNet can even find better local solutions than traditional solvers, at significantly greater speed. 
\end{itemize}

\section{Related Work}

\textbf{ML for optimization.} There has been a large body of literature leveraging ML to provide fast solutions to optimization problems, split into two major classes of approaches (see also reviews \cite{khaloie2025review, bengio2020machine, hasan2020survey}).
The first class of approaches entails using ML to learn end-to-end surrogate models,
% for a particular class of optimization problems, 
in which problem parameters that vary between optimization problem instances are treated as inputs to the ML model \cite{van2025optimization, amos2023tutorial}.
This includes supervised learning methods that require a training dataset of solved instances generated via traditional solvers \cite{chen2022learning, zamzam2020learning, fioretto2021lagrangian}, as well as self-supervised/unsupervised approaches that directly embed the optimization objective and/or constraint functions into the ML loss function \cite{donti2021dc3, tanneau2024dual, park2023compact, amos2017optnet}.
In this work, we adopt the self-supervised paradigm, which has been favored in recent work as it avoids the large amount of time and resources required to create supervised training datasets \cite{thams2019efficient}.
The second class of ML-for-optimization approaches involves using ML to improve the performance of other optimization solvers.
For instance, ML has been used to
approximate algorithm steps for convex optimization to achieve guaranteed fast convergence \cite{tan2023data}, provide warm start points to accelerate fixed-point optimization algorithms \cite{sambharya2024learning}, or approximate or remove complex or redundant constraints to reduce computational complexity \cite{bertsimas2025global, misra2022learning}.
The choice between whether to use an 
end-to-end vs.~ML-for-other-solvers approach
has historically been philosophical, with a lack of rigorous comparison. % between these approaches.
While our work is primarily situated within the end-to-end surrogates literature, we also draw inspiration from the ML-for-other-solvers approach, notably by incorporating an optimization warm-start component within our framework.

\textbf{Feasibility enforcement in ML-for-optimization.} While end-to-end ML approaches have been successful in 
% reducing the computational complexity of providing optimization solutions, 
providing fast solutions,
these approaches have struggled to satisfy problem constraints.
Several constraint-enforcement approaches have been proposed to mitigate this, 
% challenge, 
including penalty methods, projection-based methods, and iterative methods. 
Penalty methods 
% aim to prioritize 
incentivize
constraint enforcement by penalizing constraint violations within the ML loss function \cite{fioretto2021lagrangian, park2023self}; however, these methods do not have feasibility guarantees and can generate highly infeasible solutions in practice \cite{donti2021dc3,park2023self, park2023compact}. 
Projection-based methods
% , on the other hand, 
\emph{do}
provide feasibility guarantees, by projecting the output of the ML model onto the feasible set \cite{min2024hard, huang2021differentiable}---in the context of 
% NN-based methods, 
NNs,
this can be done via differentiable orthogonal projection methods that admit end-to-end training \cite{min2024hard, amos2017optnet, cvxpylayers2019}.
In special cases where the projection admits a computationally efficient special structure (e.g., 
solvable in closed form), projection methods can work extremely well, providing fast, feasible, high-quality solutions \cite{tanneau2024dual, konstantinov2023new, tordesillas2023rayen}.
However, in general cases, solving orthogonal projection problems can be expensive, making training and inference 
slow, particularly in high-dimensional problems.
Iterative methods aim to be more computationally efficient by instead deploying iterative algorithms to feasibility-correct the predictions of an ML algorithm. %, in a way that is computationally efficient and able to handle general classes of complex constraints.
Notable examples include DC3 \cite{donti2021dc3}, which uses a variable elimination-based approach to satisfy equality constraints and a gradient-based iterative approach to satisfy inequality constraints, and ENFORCE \cite{lastrucci2025enforce}, which handles equality constraints by unrolling a sequential quadratic programming procedure.
Our work presents an iterative method that handles both equality and inequality constraints and, unlike previous iterative methods, provides guarantees on constraint enforcement, which are particularly important in (e.g.) safety-critical settings.


\textbf{Differentiable optimization in deep learning.} 
Optimization problems can be directly integrated into neural networks as differentiable layers to embed relevant inductive biases or constraint satisfaction directly into the learning process. 
To enable end-to-end training, gradients must be computed through these optimization layers. 
\textit{Implicit differentiation} approaches compute the gradient implicitly using optimality conditions or fixed-point equations, with approaches presented in the literature for, e.g., quadratic programs \cite{amos2017optnet, magoon2024differentiation}, disciplined convex problems \cite{cvxpylayers2019}, cone problems \cite{donti2020enforcing, agrawal2019differentiating}, submodular minimization \cite{djolonga2017differentiable}, and other problem classes \cite{gould2021deep}. 
\textit{Unrolled differentiation}, by contrast, explicitly unrolls the optimization iterations and leverages automatic differentiation tools to track gradients through each step \cite{monga2021algorithm, kotary2023backpropagation, shaban2019truncated, scieur2022curse}.
While implicit differentiation can be more computationally efficient and numerically stable, it relies on obtaining an optimal solution to the optimization problem (which may be unavailable if the optimization did not converge) and is only applicable 
in situations where the ML-generated inputs show up in the fixed-point equations (which is not true for, e.g., warm start points).
For the latter reason, FSNet employs unrolled differentiation (with truncated gradients) 
% Our FSNet employs this technique 
to differentiate through the feasibility-seeking step during training. 
For a review of methods integrating optimization within neural networks, as well as differentiation techniques, we refer readers to \cite{mandi2024decision, blondel2022efficient}.

\section{FSNet: Feasibility-Seeking-Integrated Neural Network}

We now present FSNet, a novel NN-based framework for providing fast, feasible, high-quality solutions to constrained optimization problems. Specifically, we consider the task of solving parametric optimization problems of the form:  
\begin{align} \label{prob: nonlinear problem}
	\min_{ y \in \R^n} \;\; f(y; \x) \;\; \st \;\; g (y; \x) \leq 0, \; h (y; \x) = 0,
\end{align}
%
where $y \in \R^n$ are decision variables, $\x \in \R^d$ are parameters, $f: \R^n \times \R^d \rightarrow \R$ is the objective function, and $g:\R^n \times \R^d \rightarrow \R^{n_\mathrm{ineq}}$, 
$h: \R^n \times \R^d \rightarrow \R^{n_\mathrm{eq}}$  are inequality and equality constraint functions, respectively.
For illustrative examples of how $x$ and $y$ are defined in various practical problems, we refer the reader to the tutorial on amortized optimization \cite{amos2023tutorial}.
Since the minimizer(s) of problem \eqref{prob: nonlinear problem} change as the parameters $\x$ vary, this implies an underlying mapping (or mappings) from the parameters $\x$ to the minimizer(s).
Thus, we can frame the problem of solving the parametric optimization problem as a learning problem in which we aim to 
use an ML model to 
approximate an underlying mapping.
Specifically, let $\hy: \R^d \rightarrow \R^n$ denote an NN-based function parameterized by $\theta \in \R^{n_\theta}$, with $\hy(\x)$ being its prediction for optimization parameters $\x$. 
We aim to find $\theta$ by solving: 
\begin{align}\label{prob: learning}
    \min_{\theta \in \R^{n_\theta}} \;\; \E_{\x \sim \mathcal{D}} [f(\hy(\x); \x)] \;\; \st \;\; g (\hy(\x); \x) \leq 0, \; h (\hy(\x); \x) = 0 \;\; \forall~x \sim \mathcal{D},
\end{align}
where $\x$ follows an unknown distribution $\cD$, and we assume access only to a finite dataset sampled from this distribution.
Importantly, the goal is to choose $\theta$ to minimize the expected objective value while \emph{ensuring that predictions satisfy all constraints}.
However, due to sources of error such as optimization error, representation error, and generalization error, NNs alone often fail to produce feasible solutions. 
% Therefore, relying solely on a NN is insufficient for solving problem \eqref{prob: learning}. 
To address this, it is necessary to have an additional module that can strictly enforce the feasibility of the predicted minimizer for any $\x \sim \cD$.


To address this problem, we introduce the FSNet framework, which integrates a feasibility-seeking step into the NN training and inference pipelines. 
This framework is illustrated in Figure~\ref{fig: method diagram}, with corresponding pseudocode in Algorithm \ref{alg: method}.
FSNet consists of three main parts: \textcolor[HTML]{335b9c}{\textbf{prediction}}, \textcolor[HTML]{b04b17}{\textbf{feasibility seeking}}, and \textbf{projection-inspired loss}. 
Specifically, the prediction step entails employing any standard NN $\y: \R^d \rightarrow \R^n$ to map from input parameters $\x$ to a candidate minimizer $\y(\x)$.
As this candidate generally violates the constraints, we subsequently execute a feasibility-seeking procedure, which solves an infeasibility-minimization problem to transform $\y(\x)$ into a feasible solution $\hy(\x)$.
We train this prediction + feasibility-seeking pipeline end-to-end to minimize a loss function that penalizes both the original optimization objective $f$ and the distance between the original prediction $\y(\x)$ and the post-feasibility-seeking outputs $\hy(\x)$.
% We now describe the feasibility-seeking step and loss function in more detail.

\begin{figure}[t]
    \centering
    \begin{minipage}[t]{0.52\textwidth}
        \vspace{20pt}%
        \includegraphics[width=\linewidth]{Figures/Diagram.pdf}
        % \label{fig: method diagram}
    \end{minipage}
    \hfill
    \begin{minipage}[t]{0.47\textwidth}
        \vspace{10pt}%
      \centering
      \begin{algorithm}[H]
        \footnotesize
        \caption{FSNet Training Algorithm}
        \label{alg: method}
            \algrenewcommand\algorithmicindent{0.7em}%
            \begin{algorithmic}[1]
                \State \textbf{init.}~NN weights $\theta$, learning rate
                \Repeat
                    \State \textbf{sample} 
                    $x \sim \mathcal{D}$
                    \State \textbf{predict} $\y(\x)$ via NN
                    \State \textbf{compute} $\hy(\x) = \FS(\y(\x); \x)$ via~\eqref{eq: FS_GD}
                    \State \textbf{update} $\theta$ with $\nabla_\theta F(\y(\x), \hy(\x))$ via \eqref{eq: F}
                \Until{convergence}
            \end{algorithmic}
        \end{algorithm}
    \end{minipage}
    \vspace{-15pt}
    \caption{A schematic of FSNet, our framework for solving constrained optimization problems.}
    \label{fig: method diagram}
\end{figure}

\subsection{Feasibility-seeking step}\label{sec: fs}
We now discuss the idea and implementation of the feasibility-seeking step, which serves as the core mechanism in our framework for providing provable feasibility guarantees.
Given a potentially infeasible NN prediction $\y(\x)$, one natural idea to ensure feasibility is to establish a process that 
starts at $\y(\x)$ and iteratively moves the point toward the feasible set.
As this process progresses, the constraint violations are expected to gradually decrease and eventually vanish when a feasible point is reached.
This idea can be formalized by leveraging $\y(\x)$ as a warm start point for the following minimization problem, and letting $\hy(\x)$ be the problem's obtained solution:
\begin{align}\label{obj: fs}
   \min_{s \in \R^n} \; \phi (s;\x) := \|h(s; \x)\|_2^2 + \| g^+(s;\x) \|_2^2,
\end{align}
where $\phi(s;\x)$ quantifies the total violation of equality and inequality constraints for a given prediction $s$ and parameter $\x$, and $g^+(s;\x) = \max(g (s;\x), 0)$ is applied element-wise.
In practice, the violation-minimization problem \eqref{obj: fs} can be solved using one of many  iterative optimization methods with local and/or global convergence guarantees \cite{nocedal1999numerical}.
For the purposes of our later theoretical analysis,
we present the vanilla gradient descent method here, 
which minimizes $\phi(s;\x)$ over iterations $k$ via:
\begin{align}\label{eq: FS_GD}
        s_0 = \y(\x),  \;\;\; % \\
        s_{k+1} = s_k - \eta_\phi \nabla_s \phi(s_k; \x),
\end{align}
where the iterations start from the prediction $\y(\x)$ and move in the direction of steepest descent with properly-chosen step size $\eta_\phi$.
Let $\FS: \R^n \times \R^d \mapsto \R^n$ be a mapping that maps the initial point $\y(\x)$ and parameters $x$ to the final prediction $\hy(\x)$, in particular $\hy(\x) = \FS(\y(\x); \x) \coloneqq \lim_{k \rightarrow \infty} s_k$.
In practice, we terminate our iterative procedure after a finite number of iterations until the function value $\phi(s_k;\x)$ or its gradient drops below a small, pre-specified threshold. % and let our final solution be $\hy^K(\x) = s_K$.

To enable end-to-end training through the NN prediction and feasibility-seeking steps, we apply backpropagation through the unrolled iterates of the feasibility-seeking optimization process. 
(Since the prediction $\y(\x)$ does not appear within the optimality conditions of \eqref{obj: fs}, implicit differentiation is not applicable for computing this gradient.)


Unlike projection-based methods that solve expensive constrained optimization problems, the feasibility-seeking step solves an unconstrained problem, which is typically easier and faster to solve. %, and is well-suited for parallelization on GPUs. 
Unlike DC3 \cite{donti2021dc3}, which handles equality and inequality constraints separately, the feasibility-seeking step handles both simultaneously, which facilitates establishing feasibility guarantees (Section~\ref{sec: analysis}). 
% in addition, whereas DC3 may require implementing custom implicit layers to handle equality constraints, FSNet more readily leverages the automatic differentiation capabilities available in standard ML libraries by differentiating through unrolled optimization iterates. %, making it potentially easier to implement.

\subsection{Loss function}\label{sec: loss}
We train the NN end-to-end through the feasibility-seeking procedure by minimizing an unsupervised empirical loss function over a dataset $\{x^{(1)}, \dots, x^{(S)}\}$, defined as follows:
\begin{align}\label{eq: empirical loss}
    \cL(\theta) = \frac{1}{S} \sum_{i=1}^S F \left(\y(\x^{(i)}) \, , \, \hy(\x^\i)\right),
\end{align}
where $\hy(\x) = \FS(\y(\x); \x)$ is the output of the feasibility-seeking step, and we define the function $F(\y(\x) , \hy(\x))$ as:
\begin{align}\label{eq: F}
    F(\y(\x) , \hy(\x)) = f(\hy(\x);\x) + \frac{\rho}{2} \|\y(\x)- \hy(\x)\|_2^2, %, \quad \text{with}\; \hy(\x) = \FS(\y(\x); \x),
\end{align}
% with $\hy(\x) = \FS(\y(\x); \x)$ being the output of the feasibility-seeking step and
with penalty weight $\rho > 0$. % is a penalty weight.
The term $f(\hy(\x); \x)$ promotes low objective value, while the term $\| \y(\x)- \hy(\x)\|_2^2$ encourages minimal corrections by the feasibility-seeking step, thereby incentivizing the NN to produce predictions that are closer to feasible.
From a geometric perspective, the term $\|\y(\x) - \hy(\x)\|_2^2 $ can be interpreted as a soft approximation of a projection, loosely corresponding to finding a feasible point $\hy(\x)$ that remains close to the original prediction $\y(\x)$. %, similar in spirit to projecting $\y(\x)$ onto the feasible set.  
This term also plays a critical role in the convergence analysis of our framework, which we describe in the next section.

\section{Convergence Analysis} \label{sec: analysis}

In solving problem \eqref{prob: nonlinear problem}, FSNet is expected to predict a solution that minimizes the objective function while satisfying all constraints. 
Toward showing this, we prove that the feasibility-seeking procedure finds a feasible point (if one exists) at an exponentially fast rate (Section~\ref{subsec:fs-zero}) and that the overall FSNet training procedure converges in the sense of the expected gradient norm vanishing asymptotically (Section~\ref{subsec:nn-converge}). Finally, we characterize a condition under which FSNet's prediction $\hy(\x)$ coincides with an optimal solution of the target problem (Section~\ref{subsec:minimizer convergence}).
All proofs are given in Appendix~\ref{appendix: proofs}.

\subsection{Convergence of feasibility-seeking step}
\label{subsec:fs-zero}

As the feasibility-seeking step essentially solves an unconstrained optimization problem via gradient descent, the convergence analysis of the feasibility-seeking step reduces to an analysis of the performance of gradient descent in unconstrained optimization.
We state the relevant theorem and associated standard assumption below, and refer the reader to \cite{polyak1963gradient, karimi2016linear} for full proofs.

\begin{assumption}\label{assump:phi}
For any $x$, function $\phi(\cdot \,; \x)$ is $L_\phi$-smooth,
admits a minimizer $s^\star \in \R^n$,  and satisfies the PL condition with constant $\mu_\phi$, i.e., $\|\nabla\phi(s;\x)\|^2_2 \geq 2 \mu_\phi\bigl(\phi(s;\x)-\phi(s^\star; \x) \bigr)$.
\end{assumption}

\begin{theorem}[\cite{polyak1963gradient}] \label{thm: fs convergence}
    Let $\gamma = 1 - \mu_\phi/{L_\phi}$. 
    Under Assumption~\ref{assump:phi}, the sequence $\{s_k\}$ generated by the gradient descent \eqref{eq: FS_GD} with step size $\eta_\phi = 1/L_\phi$ converges to $s^\star$ with rate
    \begin{align}
        &\| s_k - s^\star \| \leq \frac{8 \eta_\phi L_\phi^2}{\mu_\phi^2} \gamma^{k-1} (\phi(s_0; \x) - \phi(s^\star; \x)),
         \label{eq: thm1-eq1}\\
        &\phi(s_k;\x) - \phi(s^\star; \x) \leq \gamma^k (\phi(s_0;\x) - \phi(s^\star; \x)). \label{eq: thm1-eq2}
    \end{align}
\end{theorem}
From Theorem 1, we see that if the target optimization problem \eqref{prob: nonlinear problem} admits at least one feasible point for a given $\x$, i.e.,  $\phi(s^\star; \x) = 0$, then the feasibility-seeking step drives the NN's prediction $\y(\x)$ to a feasible point at an exponentially fast rate. 
Indeed, gradient descent needs at most $\bigO (\log(1/{\epsilon_\phi}))$ iterations to reduce the violation value $\phi(s_k;\x)$ below a small threshold $\epsilon_\phi>0$. 

\subsection{Convergence of NN training}
\label{subsec:nn-converge}
This section analyzes the convergence of FSNet training using stochastic gradient descent (SGD). 
Before delving into the analysis, we emphasize an important practical aspect that can impact the convergence of training:~the feasibility-seeking problem \eqref{obj: fs} is typically solved by unrolling a finite number of gradient descent steps in \eqref{eq: FS_GD} until a small pre-specified convergence tolerance is met.
To formalize this, let $\FS^K(\y(\x); \x)$ denote the feasibility-seeking mapping obtained by unrolling gradient descent in \eqref{eq: FS_GD} for $K$ iterations. 
We then define the resulting approximate solution as  $\hy^K(\x) \coloneqq \FS^K(\y(\x); \x)$. 
Under this finite unrolling scheme, we now introduce a modified loss function to reflect the approximate nature of the feasibility-seeking step:
\begin{align}\label{eq: approx loss}
    L(\theta) = \frac{1}{S} \sum_{i=1}^S F\big(\y(\x^\i), \, \hy^K(\x^\i)\big).
\end{align}
The NN training is performed using SGD with an unbiased gradient estimator $\tgrad L(\theta_t)$:
\begin{align}\label{eq: SGD}
    \theta_{t+1} = \theta_t - \eta \tgrad L(\theta_t) \quad \text{with} \quad \mathbb{E}_t[\tgrad L(\theta_t)] = \nabla L(\theta_t),
\end{align}
where $\mathbb{E}_t[\cdot]$ is the conditional expectation given the randomness up to step $t$.

We now analyze the convergence of NN training under this finite unrolling scheme. Establishing SGD convergence for the modified loss $L$ follows standard arguments \cite{shalev2014understanding}.
However, to understand how finite unrolling of the feasibility-seeking step affects optimization of the exact loss $\cL$, we must also bound the discrepancy $\| \nabla L(\theta) - \nabla \cL(\theta)\|_2$. Thus, we need to make boundedness and $L$-smoothness assumptions to control this discrepancy; otherwise, it can be arbitrarily large.
Below, we use $J_{w} \psi(w) \coloneqq \partial \psi(w)/{\partial w}$ to denote the Jacobian of vector-valued function $\psi$ with respect to $w$.

\begin{assumption}[Boundedness] \label{assump:nn-bound}
There exist constants $B_N, B_{FS}, B_F > 0$ such that for all $\theta, x$: 
    \begin{align*}
        \| J_\theta \y(\x) \|_2 \leq B_N, \quad  \|J_y \FS^K(\y(\x);\x) \|_2 \leq B_{FS}, \quad
        \|\nabla_{\hat{y}} F(\y(\x), \hy(\x)) \|_2 \leq B_F. \quad
    \end{align*}
\end{assumption}

\begin{assumption}[Smoothness] \label{assump:nn-smooth}
The function $F(y,\hat{y})$ is $L_F$-smooth with respect to $\hat{y}$, and the function $L(\theta)$ is $L_N$-smooth and bounded below by $L^{\star}$. 
The mappings $\FS$ and $\FS^K$ satisfy for $L_{FS}>0$ that
\[
\|J_y \FS(\y(\x); \x) - J_y \FS^K(\y(\x); \x) \|_2 \leq L_{FS} \| \hy(\x) - \hy^K(\x) \|_2, \quad \forall x \in \R^d.
\]
\end{assumption}

\begin{restatable}{theorem}{thmSgdConvergence}\label{thm: SGD convergence 1}
Suppose that $\mathbb{E}_t \left[ \| \tgrad L(\theta_t) \|_2^2 \right] \leq G$ for all $t$ and that the feasibility-seeking step is unrolled for $K$ steps in every forward pass. Let $\gamma$ be as defined in Theorem~\ref{thm: fs convergence}.     
    Under Assumption~\ref{assump:phi}, ~\ref{assump:nn-bound}, and ~\ref{assump:nn-smooth}, 
    after running the SGD for $T$ iterations with step size $\eta = \eta_0 / \sqrt{T}$ for some $\eta_0 > 0$, there is at least one iteration $t \in \{0, \dots, T-1\}$ satisfying 
    \begin{align}
        &\mathbb{E} \left[ \left\| \nabla L(\theta_t) \right\|_2^2 \right] \leq \bigO\left(T^{-1/2}\right),  \label{eq: thm2-bound1}\\
        &\mathbb{E} \left[\|\nabla \cL(\theta_t) \|_2\right] \leq 
        \bigO\left( T^{-1/4} + \gamma^K \right). \label{eq: thm2-bound2}
    \end{align}
\end{restatable}

This theorem shows that, when the finite unrolling is done in practice, SGD converges in expectation to a stationary point of the NN parameters. 
Moreover, the finite unrolling introduces a bias of order $\bigO(\gamma^K)$ to the decay rate of $\nabla \cL$, but the bias decays exponentially in $K$ due to $\gamma < 1$. 
In other words, for a moderate number of unrolled steps $K$, the bias is negligible, and the gradient norms under finite and infinite unrolling diminish at virtually the same rate. 
Consequently, finite unrolling is not only computationally efficient but also sufficient to preserve convergence guarantees. 


\subsection{Convergence to minimizers}\label{subsec:minimizer convergence}
In the previous section, we analyzed parameter‐level convergence of the FSNet training process. Here, we turn to output‐level convergence: first, we derive conditions under which the network’s prediction is a stationary point of the function $F$, and then we identify when that prediction exactly matches the minimizer of the target optimization problem \eqref{prob: nonlinear problem}.

By the chain rule, the gradient of the loss~\eqref{eq: empirical loss} is the mean of $(J_\theta \y(\x^\i))^\top \nabla_y F \left(\y(\x^{(i)}) , \hy(\x^\i)\right)$ over all samples. This implies that SGD can stall for two reasons: either the gradient $\nabla_y F \left(\y(\x^{(i)}), \hy(\x^\i)\right)$ vanishes, or it remains nonzero but lies in the nullspace of $J_\theta \y(\x^\i)$. 
The former means that the prediction of the NN is a stationary point of $F \left(\y(\x^{(i)}), \hy(\x^\i)\right)$, which is desirable, since only a stationary point has a chance to be a minimizer 
of the target optimization problem \eqref{prob: nonlinear problem}. 
% The latter means that the parameter updates no longer affect the output, which we do not want because we expect NN to predict the stationary point. 
Therefore, we want to rule out the latter case. One potential approach is employing a sufficiently expressive NN so that Jacobian matrices are full-rank and stationary points of $F$ can be reached for all training inputs $x$.
We formalize this observation in the following lemma. 

\begin{restatable}{lemma}{lemmaStationary}\label{lemma: stationary prediction}
     Suppose SGD converges to $\theta^\star$, where 
     % the matrix 
     $J(\theta^\star) := \big((J_\theta y_{\theta^\star}(x^{(1)}))^\top \, \dots \, (J_\theta y_{\theta^\star}(x^{(S)}))^\top \big)$ has full column rank. 
     Then, the prediction of the NN $\y(\x^\i)$ is the stationary point of $F(\y(\x^\i), \hy(\x^\i))$ with $\hy(\x^\i)=\FS(\y(\x^\i); \x^\i)$ for all $i\in\{1, \dots, S\}$. 
\end{restatable}

This lemma means that after being successfully trained, the NN in the FSNet framework can predict a stationary point $\y(\x)$ of $F(\y(\x), \hy(\x))$. 
If $F$ is convex, then the stationary point is a global minimizer of $F$. 
However, if $F$ is nonconvex, the stationary point could be a minimizer, maximizer, or saddle point.
In such a case, additional assumptions are necessary for further analysis. 
Thus, we now make a (fairly strong) assumption that the NN's prediction is a global minimizer of $F$.
However, a minimizer of $F$ may not necessarily be a minimizer of the target problem \eqref{prob: nonlinear problem}, because $F$ differs from the original objective function $f$ in including a penalty term. 
Inspired by the penalty method in constrained optimization \cite{nocedal1999numerical}, we can increase the penalty weight $\rho$ such that the difference $\| \y (\x) - \hy(\x) \|_2$ vanishes, which drives the minimizers of $F$ and problem $\eqref{prob: nonlinear problem}$ to coincide. % with each other. 


\begin{restatable}{theorem}{thmMinimizer}\label{thm: minimizer convergence}
    Given specific parameters $\x$, suppose that the NN predicts $y_\tau$ which is a global minimizer of $F(y, \hat{y}; \rho_\tau) = f(\hat{y}; x) + \frac{\rho_\tau}{2} \| y - \hat{y}\|_2^2$ with $\hat{y} = \FS(y; x)$; that the mapping $\FS$ is continuous; and that $0 < \rho_\tau < \rho_{\tau+1}, \forall \tau$ and $\rho_\tau \rightarrow \infty$. Then, every limit point $(y^\star, \hat{y}^\star)$ of the sequence $\{(y_\tau, \hat{y}_\tau)\}$ admits $y^\star = \hat{y}^\star$, and $\hat{y}^\star$ is a global minimizer of the optimization problem \eqref{prob: nonlinear problem}. 
\end{restatable}

We acknowledge that the assumption in Theorem \ref{thm: minimizer convergence} is rather stringent, but it nonetheless provides some useful insights for practical implementation. 
In particular, the penalty term $\frac{\rho}{2} \|\y(\x)- \hy(\x)\|_2^2$ in \eqref{eq: F} needs to be increasing to yield high-quality solutions.
From our experiments, we find that choosing a sufficiently large value of $\rho$ is sufficient , and the range of effective $\rho$ is wide, making it easier to select an appropriate value.


\section{Practical Efficiency Improvements}\label{sec: improvement}
While the above FSNet framework offers theoretical guarantees, we also provide practical strategies to improve training efficiency and 
stability. 
We discuss the use of truncated gradient tracking to reduce computational overhead, and an enhancement to the loss function to improve training stability. 

\textbf{Truncated unrolled differentiation.} Inspired by \cite{shaban2019truncated}, we implement truncated unrolled differentiation to improve the computational efficiency of the FSNet framework. Specifically, when computing gradient updates, we backpropagate through only the first $K'$ out of the total $K$ iterations of the feasibility-seeking procedure, rather than all $K$ iterations.
In other words, only the computation of the first $K'$ iterations is included in the computational graph for backpropagation, while the remaining iterations act as pass-through layers without gradient tracking. 
This approach significantly reduces computational graph size, memory, and the overhead of the backward pass \cite{shaban2019truncated}. 
We empirically observe negligible loss in solution quality even for modest $K'$, and in Appendix~\ref{subsec: truncated diff analysis}, we prove that, under local strong convexity of $\phi$, the Jacobian error and the truncation-induced bias in the SGD updates decay exponentially with $K'$.

\textbf{Stabilizing early-stage training.}
With randomly initialized parameters, the NN may produce large constraint violations early in training, forcing the feasibility-seeking step to use many iterations and slowing down overall training. 
To mitigate this, we introduce an additional penalty for constraint violation when it exceeds a threshold:
    \begin{align}\label{eq: modified F}
    F(\y(\x), \hy(\x)) =
    \left\{ \begin{array}{ll}
         f(\hy(\x); \x) + \frac{\rho}{2} \| \y(\x)- \hy(\x)\|_2^2 + \rho_\phi \phi(\y(\x); \x)& \text{if } \phi(\y(\x);\x) \geq Q, \\
         f(\hy(\x); \x) + \frac{\rho}{2} \| \y(\x)- \hy(\x)\|_2^2 & \text{otherwise,}
    \end{array}
    \right.
\end{align}
where the threshold $Q$ can be big, e.g., $Q=10^3$.
When constraint violations are too large ($\phi(y; \x) \geq Q$), the extra term adds an additional penalty to the loss that pushes the NN to make more feasible predictions.
Otherwise, the extra term is removed, and the loss reverts exactly to the original form. 


\section{Experiments}\label{sec: experiments}

\textbf{Problem Classes.}
We evaluate our method on three families of synthetic problems: \textit{smooth convex}, \textit{smooth nonconvex}, and \textit{nonsmooth nonconvex}. The \textit{smooth convex} problems include quadratic programs (QPs), quadratically-constrained quadratic programs (QCQPs), and second-order cone programs (SOCPs). 
\textit{Smooth nonconvex} problems are constructed by adding elementwise sine/cosine terms into the objective and constraint functions of the smooth convex problems, while the \textit{nonsmooth nonconvex} problems further add an $\ell_2$-norm regularization term to the objective function (Appendix~\ref{appendix: problem formulations}). 
Each problem has 100 decision variables subject to 50 equality and 50 inequality constraints. 
We generate a dataset of 10000 samples for each problem following the procedure in \cite{donti2021dc3, liang2024homeomorphic}, and split it into  7000 training, 1000 validation, and 2000 test samples for all NN-based methods.
% \revise{
For the smooth convex settings, we additionally consider larger-scale problems with 500 decision variables, 200 equality
constraints, and 200 inequality constraints.
We use 17000 training, 1000 validation, and 2000 test samples for this case.
% }
Finally, we consider the problem of AC optimal power flow (ACOPF), a nonconvex, NP-hard problem in electric power systems that aims to schedule power generation at least cost while satisfying grid and device constraints; for this problem, we use the formulation and dataset from \cite{klamkin2025pglearn}.


\textbf{Baselines.} We compare our FSNet against:
\begin{itemize}[leftmargin=20pt, topsep=3pt]
    \item \textbf{Solver:} A traditional solver. We use OSQP \cite{stellato2020osqp} and SCS \cite{odonoghue:21} for convex problems, and use IPOPT \cite{wachter2006implementation} for nonconvex problems.     
    \item \textbf{Reduced Solver:} For a fairer runtime comparison, we reduce the solvers’ tolerances to achieve a level of constraint satisfaction comparable to FSNet.
    \item \textbf{Penalty:} A popular approach for addressing constraints in training NNs that adds a fixed-weight penalty for constraint violations within the loss function.
    \item \textbf{Adaptive Penalty:} An improved version of the fixed penalty method that adjusts the penalty weight based on current violations, as in \cite{fioretto2021lagrangian}.
    \item \textbf{DC3:} The Deep Constraint Completion and Correction framework from \cite{donti2021dc3}. 
\end{itemize}

We also implement the \textbf{Projection} method in the convex QP setting, which orthogonally projects the NN's prediction onto the feasible set using the qpth solver \cite{amos2017optnet}.
While we attempted to implement  Projection for other problem classes via cvxpylayers \cite{cvxpylayers2019}, this introduced significant computational overhead, making training prohibitively slow.
All network, training, and solver-related hyperparameters are in Appendix \ref{appendix: hyperparameters}.
We also provide experiments analyzing FSNet's performance with varying numbers of tracked iterations in truncated unrolled differentiation (Section~\ref{sec: improvement}), different values of the penalty weight $\rho$, and additional results for DC3 using our loss function \eqref{eq: empirical loss} in Appendix~\ref{appendix: more experiments}.


\textbf{Evaluation metrics.} 
We evaluate the methods using the following metrics: constraint violations, optimality gap, and computational time. 
We compute $\|h(\hy(\x);x)\|_1$ for equality violations and $\| g^+ (\hy(\x);x) \|_1$ for inequality violations. 
Since some trusted solvers can return the global minimum for convex problems, we take their solution $y_\mathrm{solver}$ to define the relative optimality gap by:
\begin{align*}
    \mathrm{optimality \,\, gap} = \frac{f(\hy(\x); \x) - f(y_\mathrm{solver}; x)}{\vert f(y_\mathrm{solver}; x) \vert}.
\end{align*}
For computational time, we compare two metrics: Batch and Sequential computational time.
For Batch computational time, all problem instances were solved in a setting of maximal parallelization within a particular budget of compute hardware. For Sequential computational time, all instances were solved one-at-a-time in sequence. Traditional solvers were run on 1 Intel(R) Xeon(R) CPU E5-2670 v2 and all NN methods were run on 1 NVIDIA Quadro RTX 8000 GPU. All metrics are computed over 3 runs.
Results on these metrics are summarized below, with full results in Appendix~\ref{appendix: more experiments}.

\textbf{Feasibility-seeking method.} We employ L-BFGS with backtracking line search for feasibility seeking, as it nicely balances convergence rate and scalability to high-dimensional problems \cite{nocedal1999numerical}. 

\subsection{Smooth convex problems}
\begin{figure}
    \centering
    \includegraphics[width=0.8\linewidth]{Figures/convex_results.pdf}
    \caption{Test results on 2000 instances of convex problems with 100 decision variables.
    FSNet consistently obtains near-zero constraint violations and small optimality gaps across the problem classes, whereas the baseline methods incur significant constraint violations.
    }
    % \vspace{-12pt}
    \label{fig: convex_results}
\end{figure}
\begin{figure}
    \centering
    \includegraphics[width=0.92\linewidth]{Figures/combined_times.pdf}
    \caption{Computational time on 2000 instances of convex and nonconvex (NC) problems with 100 decision variables.}
    \vspace{-12pt}
    \label{fig: combined_times}
\end{figure}

For the convex settings with 100 decision variables, Figure \ref{fig: convex_results} shows that FSNet outperforms the other baselines by obtaining near-zero constraint violations, whereas other methods incur significant constraint violations. 
For DC3, equality constraints are handled well, but inequality constraints are not satisfied across all cases; 
moreover, we observe that the performance of DC3 on the inequality constraints is sensitive to the step size of correction steps. 
In contrast, our feasibility-seeking step based on L-BFGS performs robustly across problem classes without needing to change the hyperparameters. 

%optimality gap
% Regarding the optimality gap, it is noted that the gap can be negative when the objective values given by the NN are lower than the solver. 


FSNet also obtains small optimality gaps, all less than $0.2\%$ on average. The baseline methods also exhibit small optimality gaps, but their solutions violate the constraints significantly, which means these gaps are not meaningful.
Specifically, we see that when constraint violations exist, NNs can easily find an infeasible point to achieve low objective values and optimality gaps--- e.g., Penalty, Adaptive Penalty, and even DC3 attain negative optimality gaps for some instances (Figure~\ref{fig: convex_results}). 


Overall, FSNet achieves order-of-magnitude speedups over traditional solvers---e.g., 30-1000$\times$ faster in batch solving (Figure~\ref{fig: combined_times}).
% With batch solving, this method can process 2000 problem instances in less than 0.5 seconds. 
For QP problems, FSNet outperforms OSQP when using batch solving but is slower under sequential solving; this difference arises because OSQP is specifically designed to exploit the problem structure of QPs, whereas FSNet does not yet incorporate such problem-specific optimizations.
Despite constraint satisfaction, the projection method runs about 3.6$\times$ slower than the traditional solver, and about 113$\times$ slower than FSNet in the batch setting (Figure~\ref{fig: combined_times}).
In contrast, for QCQP and SOCP problems, FSNet consistently outperforms traditional solvers in both batch and sequential settings.
Due to the incorporation of the feasibility-seeking step, FSNet has a slower runtime compared to other NN-based methods; however, this extra overhead is a necessary tradeoff to ensure the feasibility of the prediction.
Conversely, faster methods fail to ensure feasible solutions, rendering their speed meaningless in applications where constraint satisfaction is non-negotiable.


For larger-scale convex problems (with 500 variables, 200 equality constraints, and 200 inequality constraints), the results are presented in Table~\ref{tab: larger scale results}. In the batch setting, FSNet achieves speedups of 188$\times$ (QP) and 147$\times$ (QCQP), which are notably higher than the speedups obtained in the smaller problems with 100 decision variables (namely, 31$\times$ and 104$\times$, respectively). 
For SOCP, FSNet achieves a smaller speedup of 36$\times$ compared to 1125$\times$ in the 100-variable case. 
In the sequential setting, FSNet achieved speedups of 8$\times$ (QP),143$\times$ (QCQP), and 2754$\times$ (SOCP), compared to 0.8$\times$, 6$\times$, and 18$\times$ in the corresponding 100-variable cases. In general, FSNet becomes increasingly advantageous in larger-scale settings. This can be explained by the fact that the feasibility-seeking step, which is an unconstrained feasibility problem and the dominant component of FSNet, is inherently more scalable than solving the original constrained optimization problems.
% This trend is consistent with our findings on ACOPF problems and supports the observation that FSNet becomes increasingly advantageous in larger-scale settings.

\begin{table*}[t]
\centering
\scriptsize
\setlength{\tabcolsep}{5pt}
\begin{tabular}{c c c cc c}
\toprule
\multirow{2}{*}{\textbf{Method}}
& \multicolumn{1}{c}{\textbf{Equality Vio.}}
& \multicolumn{1}{c}{\textbf{Inequality Vio.}} 
& \multicolumn{2}{c}{\textbf{Optimality Gap ($\%$)}} 
& \multicolumn{1}{c}{\textbf{Runtime (s)}} \\
& Mean (Max) 
& Mean (Max)
& Mean  &Min (Max)
& Batch (sequential)\\

\midrule
\multicolumn{6}{c}{\textbf{Convex QP: \(n=500, n_{\mathrm{eq}}=200, n_{\mathrm{ineq}}=200\)}} \\
\midrule
\multirow{2}{*}{Solver} & 1.9\e{-9}\std{2.9\e{-10}} 
& 1.4\e{-6}\std{9.8\e{-8}} 
& \multirow{2}{*}{--}
& \multirow{2}{*}{--} 
& 112.962\std{0.378} \\
& (1.2\e{-6}\std{9.8\e{-8}}) 
& (9.3\e{-4}\std{1.3\e{-4}})
&&  
& (1394.024\std{94.868}) \\

\thinrule
\rowcolor{darkblue!4}
& 1.0\e{-4}\std{2.5\e{-6}} 
& 4.5\e{-6}\std{1.5\e{-7}} 
&&0.041\std{6.6\e{-4}} 
&0.599\std{1.8\e{-3}}  \\
\rowcolor{darkblue!4}
\multirow{-2}{*}{FSNet} & (4.4\e{-3}\std{8.2\e{-4}})
&(4.2\e{-4}\std{9.0\e{-5}})
&\multirow{-2}{*}{0.066\std{2.7\e{-4}}}  
&(0.119\std{4.9\e{-3}})  
& (163.558\std{1.027}) \\

\midrule
\multicolumn{6}{c}{\textbf{Convex QCQP: \(n=500, n_{\mathrm{eq}}=200, n_{\mathrm{ineq}}=200\)}} \\
\midrule
\multirow{2}{*}{Solver} & 1.3\e{-5}\std{4.3\e{-7}} 
& 5.0\e{-2}\std{9.9\e{-4}}
& \multirow{2}{*}{--}
& \multirow{2}{*}{--} 
& 5508.419\std{44.220}  \\
& (3.1\e{-4}\std{7.7\e{-5}}) 
& (2.1\e{-1}\std{7.3\e{-3}})
&& 
&(33362.147\std{107.952}) \\


\thinrule
\rowcolor{darkblue!4}
 & 2.7\e{-4}\std{1.4\e{-6}}
& 9.9\e{-7}\std{4.5\e{-8}} 
& 
& 0.087\std{2.5\e{-3}} 
& 37.291\std{0.237} \\
\rowcolor{darkblue!4}
\multirow{-2}{*}{FSNet}
& (3.7\e{-4}\std{3.3\e{-7}}) 
& (1.0\e{-5}\std{4.9\e{-7}}) 
& \multirow{-2}{*}{0.143\std{1.3\e{-3}}}
&(0.328\std{0.034})  
& (232.837\std{1.229}) \\

%%%%%%%%%%%%%%%%%%%%%%%%%%%%%%%%%%%%%%%
\midrule
\multicolumn{6}{c}{\textbf{Convex SOCP: \(n=500, n_{\mathrm{eq}}=200, n_{\mathrm{ineq}}=200\)}} 
\\
\midrule
\multirow{2}{*}{Solver} &5.8\e{-13}\std{5.2\e{-16}} 
& 2.2\e{-10}\std{3.1\e{-13}} 
& \multirow{2}{*}{--}
& \multirow{2}{*}{--} 
& 523.045\std{6.284} \\
& (2.1\e{-12}\std{3.1\e{-13}}) 
& (8.1\e{-8}\std{2.7\e{-8}})
&&
&(490084.123\std{98143.538}) \\

\thinrule
\rowcolor{darkblue!4}
& 2.7\e{-4}\std{1.4\e{-6}} 
& 0.0\std{0.0}
& 
& 0.469\std{0.246}
& 14.313\std{0.073}  \\
\rowcolor{darkblue!4}
\multirow{-2}{*}{FSNet} 
& (3.7\e{-4}\std{2.1\e{-6}}) 
& (0.0\std{0.0}) 
& \multirow{-2}{*}{0.730\std{0.027}}
&(1.124\std{0.048})
& (177.890\std{1.098}) \\
\bottomrule
\end{tabular}
\caption{Results of FSNet on 2000 instances of smooth convex problems with 500 decision variables.
}
\label{tab: larger scale results}
\end{table*}


\subsection{Smooth nonconvex problems}
On smooth nonconvex problems, FSNet consistently achieves near-zero constraint violations across nonconvex problems and achieves 85-3000$\times$ batch speedups and 2.7-117$\times$ sequential speedups over traditional solvers (Figure~\ref{fig: combined_times}). While the computational time of FSNet remains relatively stable across problem types, nonconvex solvers typically require more time than their convex counterparts, resulting in even greater speedup gains for FSNet in nonconvex settings.
As in the convex cases, the Penalty and Adaptive Penalty show poor performance on these problems with respect to constraint satisfaction. Results for DC3 are not shown, as this method sometimes diverged in the presence of nonconvex constraints due to the divergence of Newton's method in the completion procedure.


Notably, FSNet actually outperforms the traditional solvers in terms of optimality for nonconvex problems. 
Specifically, for the Nonconvex QCQP and Nonconvex SOCP problem variants, the optimality gaps are negative (Table~\ref{tab: nonconvex results}), indicating that the solutions predicted by FSNet yield lower objective values than those obtained by the solvers.  
The difference is clearly illustrated in Figure~\ref{fig: nonconvex distribution results}, where the distribution of FSNet objective values is sharper than that of the solver in Nonconvex QCQP, and FSNet's distribution for Nonconvex SOCP is tightly centered at lower objective values. 
For the Nonconvex QP, the two distributions are almost perfectly aligned, 
% and the optimality gap is also minimal,
suggesting that both methods achieve similar solution quality. 
% \revise{For full results on convex and nonconvex problems, please see Tables~\ref{tab: convex results} and  \ref{tab: nonconvex results}. 
% in Appendix~\ref{appendix: more experiments}
% }

% \revise{
\subsection{Nonsmooth nonconvex problems and AC optimal power flow}
Results on nonsmooth nonconvex problems and ACOPF problems are in Appendix~\ref{appendix: more experiments}.
% In Appendix~\ref{appendix: more experiments}, we also present results on nonsmooth nonconvex problems and AC optimal power flow problems, where 
The findings are generally consistent with those for smooth nonconvex problems.
In particular, FSNet achieves near-zero constraint violations across these problems. For nonsmooth nonconvex problems, FSNet achieves 152-3316$\times$ speedups in the batch setting and 7-80$\times$ speedups in the sequential setting compared to the traditional solver. We similarly observe that FSNet can find better local solutions with lower objective values than the traditional solver in the nonsmooth nonconvex setting. 
For the ACOPF problem, FSNet is faster than IPOPT, achieving 3$\times$ speedups on the IEEE 30-bus system and 11$\times$ on the IEEE 118-bus system in the sequential computation, with a less than 1\% optimality gap. % and satisfying constraints. 
% }


\begin{figure}[t]
    \centering
    \includegraphics[width=0.7\linewidth]{Figures/histograms.pdf}
    \caption{Distributions of objective values for solvers and FSNet in nonconvex problems. FSNet obtains sharper and lower-centered objective distributions for Nonconvex QCQP and Nonconvex SOCP, while matching the solver's optimality in Nonconvex QP with nearly identical distributions.}
    \vspace{-12pt}
    \label{fig: nonconvex distribution results}
\end{figure}
% \begin{figure}
%     \centering
%     \includegraphics[width=0.75\linewidth]{Figures/nonconvex_times.pdf}
%     \caption{Computational time on 2000 instances of nonconvex problems.}
%     \label{fig: combined_times}
% \end{figure}





\section{Conclusion}\label{sec: conclusion}
We proposed a novel framework, FSNet, that integrates a feasibility-seeking step into NN training to solve general constrained optimization problems, with theoretical guarantees on feasibility and convergence.
The feasibility-seeking step solves an unconstrained optimization over a function of constraint violations, initialized at the NN's prediction, to produce a feasible solution. 
Our experiments on smooth/nonsmooth and convex/nonconvex problems demonstrate that FSNet consistently outperforms popular baseline methods in achieving negligible constraint violations. 
Regarding optimality, FSNet matches the performance of traditional solvers on convex problems and often finds better local solutions in nonconvex cases, all while being significantly faster.
% By combining accuracy, feasibility, and theoretical guarantees, FSNet makes NN-based optimization more reliable and trustworthy for real-world applications that require real-time solutions. 
% enabling real-time solutions to problems that were previously expensive to solve. 

We observe through the experiments that FSNet can perform well for problems that violate the assumptions used in the convergence analysis (e.g., $L$-smoothness, PL condition). Thus, extending the convergence analysis in Section~\ref{sec: analysis} to cover such cases is a potential direction for future work.
While FSNet shows strong performance across diverse problem classes, scalability to extremely large-scale settings also remains an open direction. 
Another potential future direction is further scaling up the performance of the feasibility-seeking step using more efficient optimization algorithms.

\clearpage
\section*{Acknowledgments}
This work was supported by the U.S. National Science Foundation under award \#2325956.

% \section*{References}
\medskip
{
\small
\bibliography{references}
\bibliographystyle{unsrt}
}

%%%%%%%%%%%%%%%%%%%%%%%%%%%%%%%%%%%%%%%%%%%%%%%%%%%%%%%%%%%%
\newpage
\appendix
\numberwithin{equation}{section}
\numberwithin{figure}{section}
\numberwithin{table}{section}

% \clearpage
% \section*{NeurIPS Paper Checklist}
% \begin{enumerate}

% \item {\bf Claims}
%     \item[] Question: Do the main claims made in the abstract and introduction accurately reflect the paper's contributions and scope?
%     \item[] Answer: \answerYes{} % Replace by \answerYes{}, \answerNo{}, or \answerNA{}.
%     \item[] Justification: The abstract and introduction clearly state the paper's focus on the design of a new framework for solving constrained optimization using neural networks. The contributions are accurately reflected, and the key assumptions for the theoretical guarantees are explicitly stated in the contribution summary.
%     \item[] Guidelines:
%     \begin{itemize}
%         \item The answer NA means that the abstract and introduction do not include the claims made in the paper.
%         \item The abstract and/or introduction should clearly state the claims made, including the contributions made in the paper and important assumptions and limitations. A No or NA answer to this question will not be perceived well by the reviewers. 
%         \item The claims made should match theoretical and experimental results, and reflect how much the results can be expected to generalize to other settings. 
%         \item It is fine to include aspirational goals as motivation as long as it is clear that these goals are not attained by the paper. 
%     \end{itemize}

% \item {\bf Limitations}
%     \item[] Question: Does the paper discuss the limitations of the work performed by the authors?
%     \item[] Answer: \answerYes{} % Replace by \answerYes{}, \answerNo{}, or \answerNA{}.
%     \item[] Justification: We discuss the strength of assumptions in Section~\ref{sec: analysis}, and the limitations of the proposed method in Section~\ref{sec: conclusion}.
%     \item[] Guidelines:
%     \begin{itemize}
%         \item The answer NA means that the paper has no limitation while the answer No means that the paper has limitations, but those are not discussed in the paper. 
%         \item The authors are encouraged to create a separate "Limitations" section in their paper.
%         \item The paper should point out any strong assumptions and how robust the results are to violations of these assumptions (e.g., independence assumptions, noiseless settings, model well-specification, asymptotic approximations only holding locally). The authors should reflect on how these assumptions might be violated in practice and what the implications would be.
%         \item The authors should reflect on the scope of the claims made, e.g., if the approach was only tested on a few datasets or with a few runs. In general, empirical results often depend on implicit assumptions, which should be articulated.
%         \item The authors should reflect on the factors that influence the performance of the approach. For example, a facial recognition algorithm may perform poorly when image resolution is low or images are taken in low lighting. Or a speech-to-text system might not be used reliably to provide closed captions for online lectures because it fails to handle technical jargon.
%         \item The authors should discuss the computational efficiency of the proposed algorithms and how they scale with dataset size.
%         \item If applicable, the authors should discuss possible limitations of their approach to address problems of privacy and fairness.
%         \item While the authors might fear that complete honesty about limitations might be used by reviewers as grounds for rejection, a worse outcome might be that reviewers discover limitations that aren't acknowledged in the paper. The authors should use their best judgment and recognize that individual actions in favor of transparency play an important role in developing norms that preserve the integrity of the community. Reviewers will be specifically instructed to not penalize honesty concerning limitations.
%     \end{itemize}

% \item {\bf Theory assumptions and proofs}
%     \item[] Question: For each theoretical result, does the paper provide the full set of assumptions and a complete (and correct) proof?
%     \item[] Answer: \answerYes{} % Replace by \answerYes{}, \answerNo{}, or \answerNA{}.
%     \item[] Justification: We provide all assumptions and proofs for our theoretical results in Section~\ref{sec: analysis} and Appendix~\ref{appendix: proofs}.
%     \item[] Guidelines:
%     \begin{itemize}
%         \item The answer NA means that the paper does not include theoretical results. 
%         \item All the theorems, formulas, and proofs in the paper should be numbered and cross-referenced.
%         \item All assumptions should be clearly stated or referenced in the statement of any theorems.
%         \item The proofs can either appear in the main paper or the supplemental material, but if they appear in the supplemental material, the authors are encouraged to provide a short proof sketch to provide intuition. 
%         \item Inversely, any informal proof provided in the core of the paper should be complemented by formal proofs provided in appendix or supplemental material.
%         \item Theorems and Lemmas that the proof relies upon should be properly referenced. 
%     \end{itemize}

%     \item {\bf Experimental result reproducibility}
%     \item[] Question: Does the paper fully disclose all the information needed to reproduce the main experimental results of the paper to the extent that it affects the main claims and/or conclusions of the paper (regardless of whether the code and data are provided or not)?
%     \item[] Answer: \answerYes{} % Replace by \answerYes{}, \answerNo{}, or \answerNA{}.
%     \item[] Justification: We report all experiment settings in Section~\ref{sec: experiments} and Appendix~\ref{appendix: hyperparameters}.
%     \item[] Guidelines:
%     \begin{itemize}
%         \item The answer NA means that the paper does not include experiments.
%         \item If the paper includes experiments, a No answer to this question will not be perceived well by the reviewers: Making the paper reproducible is important, regardless of whether the code and data are provided or not.
%         \item If the contribution is a dataset and/or model, the authors should describe the steps taken to make their results reproducible or verifiable. 
%         \item Depending on the contribution, reproducibility can be accomplished in various ways. For example, if the contribution is a novel architecture, describing the architecture fully might suffice, or if the contribution is a specific model and empirical evaluation, it may be necessary to either make it possible for others to replicate the model with the same dataset, or provide access to the model. In general. releasing code and data is often one good way to accomplish this, but reproducibility can also be provided via detailed instructions for how to replicate the results, access to a hosted model (e.g., in the case of a large language model), releasing of a model checkpoint, or other means that are appropriate to the research performed.
%         \item While NeurIPS does not require releasing code, the conference does require all submissions to provide some reasonable avenue for reproducibility, which may depend on the nature of the contribution. For example
%         \begin{enumerate}
%             \item If the contribution is primarily a new algorithm, the paper should make it clear how to reproduce that algorithm.
%             \item If the contribution is primarily a new model architecture, the paper should describe the architecture clearly and fully.
%             \item If the contribution is a new model (e.g., a large language model), then there should either be a way to access this model for reproducing the results or a way to reproduce the model (e.g., with an open-source dataset or instructions for how to construct the dataset).
%             \item We recognize that reproducibility may be tricky in some cases, in which case authors are welcome to describe the particular way they provide for reproducibility. In the case of closed-source models, it may be that access to the model is limited in some way (e.g., to registered users), but it should be possible for other researchers to have some path to reproducing or verifying the results.
%         \end{enumerate}
%     \end{itemize}


% \item {\bf Open access to data and code}
%     \item[] Question: Does the paper provide open access to the data and code, with sufficient instructions to faithfully reproduce the main experimental results, as described in supplemental material?
%     \item[] Answer: \answerYes{} % Replace by \answerYes{}, \answerNo{}, or \answerNA{}.
%     \item[] Justification: We provide the anonymized code and data with instructions to reproduce our experiments.
%     \item[] Guidelines:
%     \begin{itemize}
%         \item The answer NA means that paper does not include experiments requiring code.
%         \item Please see the NeurIPS code and data submission guidelines (\url{https://nips.cc/public/guides/CodeSubmissionPolicy}) for more details.
%         \item While we encourage the release of code and data, we understand that this might not be possible, so “No” is an acceptable answer. Papers cannot be rejected simply for not including code, unless this is central to the contribution (e.g., for a new open-source benchmark).
%         \item The instructions should contain the exact command and environment needed to run to reproduce the results. See the NeurIPS code and data submission guidelines (\url{https://nips.cc/public/guides/CodeSubmissionPolicy}) for more details.
%         \item The authors should provide instructions on data access and preparation, including how to access the raw data, preprocessed data, intermediate data, and generated data, etc.
%         \item The authors should provide scripts to reproduce all experimental results for the new proposed method and baselines. If only a subset of experiments are reproducible, they should state which ones are omitted from the script and why.
%         \item At submission time, to preserve anonymity, the authors should release anonymized versions (if applicable).
%         \item Providing as much information as possible in supplemental material (appended to the paper) is recommended, but including URLs to data and code is permitted.
%     \end{itemize}


% \item {\bf Experimental setting/details}
%     \item[] Question: Does the paper specify all the training and test details (e.g., data splits, hyperparameters, how they were chosen, type of optimizer, etc.) necessary to understand the results?
%     \item[] Answer: \answerYes{} % Replace by \answerYes{}, \answerNo{}, or \answerNA{}.
%     \item[] Justification: We provide all hyperparameters of methods in Section~\ref{sec: experiments} and Appendix~\ref{appendix: hyperparameters}.
%     \item[] Guidelines:
%     \begin{itemize}
%         \item The answer NA means that the paper does not include experiments.
%         \item The experimental setting should be presented in the core of the paper to a level of detail that is necessary to appreciate the results and make sense of them.
%         \item The full details can be provided either with the code, in appendix, or as supplemental material.
%     \end{itemize}

% \item {\bf Experiment statistical significance}
%     \item[] Question: Does the paper report error bars suitably and correctly defined or other appropriate information about the statistical significance of the experiments?
%     \item[] Answer: \answerYes{} % Replace by \answerYes{}, \answerNo{}, or \answerNA{}.
%     \item[] Justification: We run our experiments with three independent runs and report the mean and standard deviation of evaluation metrics (Section \ref{sec: experiments}). 
%     \item[] Guidelines:
%     \begin{itemize}
%         \item The answer NA means that the paper does not include experiments.
%         \item The authors should answer "Yes" if the results are accompanied by error bars, confidence intervals, or statistical significance tests, at least for the experiments that support the main claims of the paper.
%         \item The factors of variability that the error bars are capturing should be clearly stated (for example, train/test split, initialization, random drawing of some parameter, or overall run with given experimental conditions).
%         \item The method for calculating the error bars should be explained (closed form formula, call to a library function, bootstrap, etc.)
%         \item The assumptions made should be given (e.g., Normally distributed errors).
%         \item It should be clear whether the error bar is the standard deviation or the standard error of the mean.
%         \item It is OK to report 1-sigma error bars, but one should state it. The authors should preferably report a 2-sigma error bar than state that they have a 96\% CI, if the hypothesis of Normality of errors is not verified.
%         \item For asymmetric distributions, the authors should be careful not to show in tables or figures symmetric error bars that would yield results that are out of range (e.g. negative error rates).
%         \item If error bars are reported in tables or plots, The authors should explain in the text how they were calculated and reference the corresponding figures or tables in the text.
%     \end{itemize}

% \item {\bf Experiments compute resources}
%     \item[] Question: For each experiment, does the paper provide sufficient information on the computer resources (type of compute workers, memory, time of execution) needed to reproduce the experiments?
%     \item[] Answer: \answerYes{} % Replace by \answerYes{}, \answerNo{}, or \answerNA{}.
%     \item[] Justification: We provide the information on the CPU and GPU we used for our experiments in Section \ref{sec: experiments}.
%     \item[] Guidelines:
%     \begin{itemize}
%         \item The answer NA means that the paper does not include experiments.
%         \item The paper should indicate the type of compute workers CPU or GPU, internal cluster, or cloud provider, including relevant memory and storage.
%         \item The paper should provide the amount of compute required for each of the individual experimental runs as well as estimate the total compute. 
%         \item The paper should disclose whether the full research project required more compute than the experiments reported in the paper (e.g., preliminary or failed experiments that didn't make it into the paper). 
%     \end{itemize}
    
% \item {\bf Code of ethics}
%     \item[] Question: Does the research conducted in the paper conform, in every respect, with the NeurIPS Code of Ethics \url{https://neurips.cc/public/EthicsGuidelines}?
%     \item[] Answer: \answerYes{} % Replace by \answerYes{}, \answerNo{}, or \answerNA{}.
%     \item[] Justification: Our research adheres to every aspect of the NeurIPS Code of Ethics.
%     \item[] Guidelines: 
%     \begin{itemize}
%         \item The answer NA means that the authors have not reviewed the NeurIPS Code of Ethics.
%         \item If the authors answer No, they should explain the special circumstances that require a deviation from the Code of Ethics.
%         \item The authors should make sure to preserve anonymity (e.g., if there is a special consideration due to laws or regulations in their jurisdiction).
%     \end{itemize}


% \item {\bf Broader impacts}
%     \item[] Question: Does the paper discuss both potential positive societal impacts and negative societal impacts of the work performed?
%     \item[] Answer: \answerNA{} % Replace by \answerYes{}, \answerNo{}, or \answerNA{}.
%     \item[] Justification: This work focuses on methodological advances for constrained optimization and does not involve direct deployment or societal-facing applications. 
%     \item[] Guidelines:
%     \begin{itemize}
%         \item The answer NA means that there is no societal impact of the work performed.
%         \item If the authors answer NA or No, they should explain why their work has no societal impact or why the paper does not address societal impact.
%         \item Examples of negative societal impacts include potential malicious or unintended uses (e.g., disinformation, generating fake profiles, surveillance), fairness considerations (e.g., deployment of technologies that could make decisions that unfairly impact specific groups), privacy considerations, and security considerations.
%         \item The conference expects that many papers will be foundational research and not tied to particular applications, let alone deployments. However, if there is a direct path to any negative applications, the authors should point it out. For example, it is legitimate to point out that an improvement in the quality of generative models could be used to generate deepfakes for disinformation. On the other hand, it is not needed to point out that a generic algorithm for optimizing neural networks could enable people to train models that generate Deepfakes faster.
%         \item The authors should consider possible harms that could arise when the technology is being used as intended and functioning correctly, harms that could arise when the technology is being used as intended but gives incorrect results, and harms following from (intentional or unintentional) misuse of the technology.
%         \item If there are negative societal impacts, the authors could also discuss possible mitigation strategies (e.g., gated release of models, providing defenses in addition to attacks, mechanisms for monitoring misuse, mechanisms to monitor how a system learns from feedback over time, improving the efficiency and accessibility of ML).
%     \end{itemize}
    
% \item {\bf Safeguards}
%     \item[] Question: Does the paper describe safeguards that have been put in place for responsible release of data or models that have a high risk for misuse (e.g., pretrained language models, image generators, or scraped datasets)?
%     \item[] Answer: \answerNA{} % Replace by \answerYes{}, \answerNo{}, or \answerNA{}.
%     \item[] Justification: This study does not involve data or models that pose a risk of misuse. 
%     \item[] Guidelines:
%     \begin{itemize}
%         \item The answer NA means that the paper poses no such risks.
%         \item Released models that have a high risk for misuse or dual-use should be released with necessary safeguards to allow for controlled use of the model, for example by requiring that users adhere to usage guidelines or restrictions to access the model or implementing safety filters. 
%         \item Datasets that have been scraped from the Internet could pose safety risks. The authors should describe how they avoided releasing unsafe images.
%         \item We recognize that providing effective safeguards is challenging, and many papers do not require this, but we encourage authors to take this into account and make a best faith effort.
%     \end{itemize}

% \item {\bf Licenses for existing assets}
%     \item[] Question: Are the creators or original owners of assets (e.g., code, data, models), used in the paper, properly credited and are the license and terms of use explicitly mentioned and properly respected?
%     \item[] Answer: \answerYes{} % Replace by \answerYes{}, \answerNo{}, or \answerNA{}.
%     \item[] Justification: We cite and credit all used assets.
%     \item[] Guidelines:
%     \begin{itemize}
%         \item The answer NA means that the paper does not use existing assets.
%         \item The authors should cite the original paper that produced the code package or dataset.
%         \item The authors should state which version of the asset is used and, if possible, include a URL.
%         \item The name of the license (e.g., CC-BY 4.0) should be included for each asset.
%         \item For scraped data from a particular source (e.g., website), the copyright and terms of service of that source should be provided.
%         \item If assets are released, the license, copyright information, and terms of use in the package should be provided. For popular datasets, \url{paperswithcode.com/datasets} has curated licenses for some datasets. Their licensing guide can help determine the license of a dataset.
%         \item For existing datasets that are re-packaged, both the original license and the license of the derived asset (if it has changed) should be provided.
%         \item If this information is not available online, the authors are encouraged to reach out to the asset's creators.
%     \end{itemize}

% \item {\bf New assets}
%     \item[] Question: Are new assets introduced in the paper well documented and is the documentation provided alongside the assets?
%     \item[] Answer: \answerYes{} % Replace by \answerYes{}, \answerNo{}, or \answerNA{}.
%     \item[] Justification: We provide anonymized code and (synthetic) data for the paper with instructions for running the code.
%     \item[] Guidelines:
%     \begin{itemize}
%         \item The answer NA means that the paper does not release new assets.
%         \item Researchers should communicate the details of the dataset/code/model as part of their submissions via structured templates. This includes details about training, license, limitations, etc. 
%         \item The paper should discuss whether and how consent was obtained from people whose asset is used.
%         \item At submission time, remember to anonymize your assets (if applicable). You can either create an anonymized URL or include an anonymized zip file.
%     \end{itemize}

% \item {\bf Crowdsourcing and research with human subjects}
%     \item[] Question: For crowdsourcing experiments and research with human subjects, does the paper include the full text of instructions given to participants and screenshots, if applicable, as well as details about compensation (if any)? 
%     \item[] Answer: \answerNA{} % Replace by \answerYes{}, \answerNo{}, or \answerNA{}.
%     \item[] Justification: We do not conduct crowdsourcing experiments nor research with human subjects.
%     \item[] Guidelines:
%     \begin{itemize}
%         \item The answer NA means that the paper does not involve crowdsourcing nor research with human subjects.
%         \item Including this information in the supplemental material is fine, but if the main contribution of the paper involves human subjects, then as much detail as possible should be included in the main paper. 
%         \item According to the NeurIPS Code of Ethics, workers involved in data collection, curation, or other labor should be paid at least the minimum wage in the country of the data collector. 
%     \end{itemize}

% \item {\bf Institutional review board (IRB) approvals or equivalent for research with human subjects}
%     \item[] Question: Does the paper describe potential risks incurred by study participants, whether such risks were disclosed to the subjects, and whether Institutional Review Board (IRB) approvals (or an equivalent approval/review based on the requirements of your country or institution) were obtained?
%     \item[] Answer: \answerNA{} % Replace by \answerYes{}, \answerNo{}, or \answerNA{}.
%     \item[] Justification: This study does not involve crowdsourcing nor research with human subjects. 
%     \item[] Guidelines:
%     \begin{itemize}
%         \item The answer NA means that the paper does not involve crowdsourcing nor research with human subjects.
%         \item Depending on the country in which research is conducted, IRB approval (or equivalent) may be required for any human subjects research. If you obtained IRB approval, you should clearly state this in the paper. 
%         \item We recognize that the procedures for this may vary significantly between institutions and locations, and we expect authors to adhere to the NeurIPS Code of Ethics and the guidelines for their institution. 
%         \item For initial submissions, do not include any information that would break anonymity (if applicable), such as the institution conducting the review.
%     \end{itemize}

% \item {\bf Declaration of LLM usage}
%     \item[] Question: Does the paper describe the usage of LLMs if it is an important, original, or non-standard component of the core methods in this research? Note that if the LLM is used only for writing, editing, or formatting purposes and does not impact the core methodology, scientific rigorousness, or originality of the research, declaration is not required.
%     %this research? 
%     \item[] Answer: \answerNA{} % Replace by \answerYes{}, \answerNo{}, or \answerNA{}.
%     \item[] Justification: This paper does not involve LLMs as any important, original, or non-standard components. 
%     \item[] Guidelines:
%     \begin{itemize}
%         \item The answer NA means that the core method development in this research does not involve LLMs as any important, original, or non-standard components.
%         \item Please refer to our LLM policy (\url{https://neurips.cc/Conferences/2025/LLM}) for what should or should not be described.
%     \end{itemize}

% \end{enumerate}


\clearpage

% \maketitle
\section{Problem formulations} \label{appendix: problem formulations}
\subsection{Convex problems}
The following formulations share: $y \in \R^n$, $Q \in \mathbb{S}^n_{++}, p \in \R^n, A \in \R^{n_\text{eq} \times n}, x \in \R^{n_\text{eq}}$, $L, U \in \R^n$.
\begin{align*}
    \textbf{QP:} \quad \min_{L \leq y \leq U}\,\, & \frac{1}{2} y^\top Q y + p^\top y \\
    \st\,\, & Ay = x, \quad Gy \leq h,
\end{align*}
where $G \in \R^{n_\text{ineq} \times n}, h \in \R^{n_\text{ineq}}$.
\begin{align*}
    \textbf{QCQP:} \quad \min_{L \leq y \leq U}\,\, & \frac{1}{2} y^\top Q y + p^\top y \\
    \st\,\, & Ay = x, \quad y^\top H_i y + g_i^\top y \leq h_i,
\end{align*}
where 
% $x \in \R^n$, $Q \in \mathbb{S}^n_{++}, q \in \R^n, A \in \R^{n_\text{eq} \times n}, p \in \R^{n_\text{eq}}$, 
$H_i \in \mathbb{S}^n_{++}, g_i \in \R^n, h \in \R$ for $i=1, \dots, n_\text{ineq}$.
\begin{align*}
    \textbf{SOCP:} \quad \min_{L \leq y \leq U}\,\, & \frac{1}{2} y^\top Q y + p^\top y \\
    \st\,\, & Ay = x, \quad \| G_i y + h_i \|_2 \leq c_i^\top y + d_i,
\end{align*}
where $G_i \in \R^{m\times n}, h_i \in \R^m, c_i \in \R^n, d_i \in \R$ for $i=1, \dots, n_\text{ineq}$.

\subsection{Nonconvex problems}
We modify the convex problems by adding element-wise sine and cosine functions. 
\begin{align*}
    \textbf{Nonconvex QP:} \quad \min_{L \leq y \leq U}\,\, & \frac{1}{2} y^\top Q y + p^\top \sin(y) \\
    \st\,\, & Ay = x, \quad G \sin(y) \leq h\cos(\x), \\
% \end{align*}
% \begin{align*}
    \textbf{Nonconvex QCQP:} \quad \min_{L \leq y \leq U}\,\, & \frac{1}{2} y^\top Q y + p^\top \sin(y) \\
    \st\,\, & Ay = x, \quad y^\top H_i y + g_i^\top \cos(y) \leq h_i,\\
% \end{align*}
% \begin{align*}
    \textbf{Nonconvex SOCP:} \quad \min_{L \leq y \leq U}\,\, & \frac{1}{2} y^\top Q y + p^\top \sin(y) \\
    \st\,\, & Ay = x, \quad \| G_i \cos(y) + h_i \|_2 \leq c_i^\top y + d_i.
\end{align*}

\subsection{Nonsmooth nonconvex problems}
Compared to the (smooth) convex problems, we further add an $\ell_2$-norm regularization term to the objective function.
\begin{align*}
    \textbf{Nonsmooth nonconvex QP:} \quad \min_{L \leq y \leq U}\,\, & \frac{1}{2} y^\top Q y + p^\top \sin(y) + \lambda \|y \|_2 \\
    \st\,\, & Ay = x, \quad G \sin(y) \leq h\cos(\x), \\
% \end{align*}
% \begin{align*}
    \textbf{Nonsmooth nonconvex QCQP:} \quad \min_{L \leq y \leq U}\,\, & \frac{1}{2} y^\top Q y + p^\top \sin(y) + \lambda \|y \|_2\\
    \st\,\, & Ay = x, \quad y^\top H_i y + g_i^\top \cos(y) \leq h_i, \\
% \end{align*}
% \begin{align*}
    \textbf{Nonsmooth nonconvex SOCP:} \quad \min_{L \leq y \leq U}\,\, & \frac{1}{2} y^\top Q y + p^\top \sin(y) + \lambda \|y \|_2\\
    \st\,\, & Ay = x, \quad \| G_i \cos(y) + h_i \|_2 \leq c_i^\top y + d_i.
\end{align*}

% \revise{
\subsection{AC optimal power flow}

For the AC optimal power flow problem, we use the formulation from \cite{klamkin2025pglearn}, reproduced below:


% \small
\begin{align*}
    \min_{p^\mathrm{g}, q^\mathrm{g}, p^\mathrm{f}, q^\mathrm{f}, p^\mathrm{t}, q^\mathrm{t}, v, \theta} \quad & \sum_{i \in \mathcal{N}} \sum_{j \in \mathcal{G}_i} c_j p_j^\mathrm{g} \\
    \text{s.t.} \quad 
    & \sum_{j \in \mathcal{G}_i} p_j^\mathrm{g} - \sum_{j \in \mathcal{L}_i} p_j^\mathrm{d} - g_i^\mathrm{s} v_i^2 = \sum_{e \in \mathcal{E}_i} p^\mathrm{f}_e + \sum_{e \in \mathcal{E}_i^R} p_e^\mathrm{t} & \forall i \in \mathcal{N} \\
    & \sum_{j \in \mathcal{G}_i} q_j^\mathrm{g} - \sum_{j \in \mathcal{L}_i} q_j^\mathrm{d} + b_i^\mathrm{s} v_i^2 = \sum_{e \in \mathcal{E}_i} q^\mathrm{f}_e + \sum_{e \in \mathcal{E}_i^R} q_e^\mathrm{t} & \forall i \in \mathcal{N} \\
    & p_e^\mathrm{f} = g_e^\mathrm{tt} v_i^2 + g_e^\mathrm{ft} v_i v_j \cos(\theta_i - \theta_j) + b_e^\mathrm{ft} v_i v_j \sin(\theta_i - \theta_j) & \forall e = (i,j) \in \mathcal{E} \\
    & q_e^\mathrm{f} = -b_e^\mathrm{tt} v_i^2 - b_e^\mathrm{ft} v_i v_j \cos(\theta_i - \theta_j) + g_e^\mathrm{ft} v_i v_j \sin(\theta_i - \theta_j) & \forall e = (i,j) \in \mathcal{E} \\
    & p_e^\mathrm{t} = g_e^\mathrm{tt} v_j^2 + g_e^\mathrm{tf} v_i v_j \cos(\theta_i - \theta_j) - b_e^\mathrm{tf} v_i v_j \sin(\theta_i - \theta_j) & \forall e = (i,j) \in \mathcal{E} \\
    & q_e^\mathrm{t} = -b_e^\mathrm{tt} v_j^2 - b_e^\mathrm{tf} v_i v_j \cos(\theta_i - \theta_j) - g_e^\mathrm{tf} v_i v_j \sin(\theta_i - \theta_j) & \forall e = (i,j) \in \mathcal{E} \\
    & (p_e^\mathrm{f})^2 + (q_e^\mathrm{f})^2 \leq \overline{S}_e^2 & \forall e \in \mathcal{E} \\
    & (p_e^\mathrm{t})^2 + (q_e^\mathrm{t})^2 \leq \overline{S}_e^2 & \forall e \in \mathcal{E} \\
    & \underline{\Delta \theta}_e \leq \theta_i - \theta_j \leq \overline{\Delta \theta}_e & \forall e = (i,j) \in \mathcal{E} \\
    & \theta^{\text{ref}} = 0 \\
    & \underline{p}_i^\mathrm{g} \leq p_i^\mathrm{g} \leq \overline{p}_i^\mathrm{g} & \forall i \in \mathcal{G} \\
    & \underline{q}_i^\mathrm{g} \leq q_i^\mathrm{g} \leq \overline{q}_i^\mathrm{g} & \forall i \in \mathcal{G} \\
    & \underline{v}_i \leq v_i \leq \overline{v}_i & \forall i \in \mathcal{N} \\
    & -\overline{S}_e \leq p_e^\mathrm{f} \leq \overline{S}_e & \forall e \in \mathcal{E} \\
    & -\overline{S}_e \leq q_e^\mathrm{f} \leq \overline{S}_e & \forall e \in \mathcal{E} \\
    & -\overline{S}_e \leq p_e^\mathrm{t} \leq \overline{S}_e & \forall e \in \mathcal{E} \\
    & -\overline{S}_e \leq q_e^\mathrm{t} \leq \overline{S}_e & \forall e \in \mathcal{E}
\end{align*}
where $\cN$, $\cE$, and $\cG$ denote the sets of buses, branches, and generators, respectively; $p^\mathrm{g}$ and $q^\mathrm{g}$ represent the active and reactive generation powers; $p^\mathrm{f}$ and $ q^\mathrm{f}$ ($p^\mathrm{t}$ and $ q^\mathrm{t}$) denote the active and reactive power flows in the forward (reverse) direction; and $v$ and $\theta$ denote the voltage magnitude and angle. The network parameters $g^\mathrm{ff}, g^\mathrm{ft}, g^\mathrm{tf}, g^\mathrm{tt}$ and $b^\mathrm{ff}, b^\mathrm{ft}, b^\mathrm{tf}, b^\mathrm{tt}$ are obtained from the admittance matrix of the system.
In addition, $\overline{S}$ denotes the apparent power flow limit, while $\underline{\Delta \theta}$ and $\overline{\Delta \theta}$ specify the lower and upper bounds on the voltage angle difference. The parameters $\underline{p}^\mathrm{g}$ and $\overline{p}^\mathrm{g}$ ($\underline{q}^\mathrm{g}$ and $\overline{q}^\mathrm{g}$) indicate the lower and upper limits of active (reactive) generation, and $\underline{v}$ and $\overline{v}$ define the permissible range of voltage magnitudes.
% }


\clearpage
\section{More experiment results}\label{appendix: more experiments}
\subsection{Smooth convex problems}
Table~\ref{tab: convex results} presents detailed results for smooth convex problems. Figure~\ref{fig: time wrt batch} plots FSNet’s runtime on 2000 test instances as a function of batch size. 
Thanks to GPU parallelism, the runtime decreases monotonically as batch size increases, so it is recommended to use the largest batch size possible to maximize speedup. 
Figure~\ref{fig: CPU usage} reports CPU utilization for the parallel solver runs. Across all problems, CPU usage remains near 100$\%$ throughout execution, showing that we fully exploit available CPU resources. This indicates that our runtime comparisons between NN-based methods and the solver are fair.

\begin{table*}[ht]
\centering
\scriptsize
\setlength{\tabcolsep}{5pt}
\begin{tabular}{ccc cc c}
\toprule
\multirow{2}{*}{\textbf{Method}}
& \multicolumn{1}{c}{\textbf{Equality Vio.}}
& \multicolumn{1}{c}{\textbf{Inequality Vio.}} 
& \multicolumn{2}{c}{\textbf{Optimality Gap ($\%$)}} 
& \multicolumn{1}{c}{\textbf{Runtime (s)}} \\
& Mean (Max) 
& Mean (Max) 
& Mean  & Min (Max)
& Batch (Sequential)\\

\midrule
\multicolumn{6}{c}{\textbf{Convex QP: \(n=100, n_{\mathrm{eq}}=100, n_{\mathrm{ineq}}=100\)}} \\
\midrule
\multirow{2}{*}{Solver} & 1.3\e{-13}\std{1.7\e{-16}} 
& 2.4\e{-14}\std{2.44\e{-16}} 
& \multirow{2}{*}{--}
& \multirow{2}{*}{--} 
& 3.802\std{0.045}\\
& (1.9\e{-13}\std{5.5\e{-15}}) 
& (5.4\e{-14}\std{4.19\e{-16}}) 
&&  
& (140.328\std{4.585}) \\
\thinrule
\multirow{2}{*}{Reduced Solver} & 1.2\e{-6}\std{7.4\e{-8}} 
& 2.1\e{-4}\std{1.1\e{-5}} 
& \multirow{2}{*}{2.5\e{-4}\std{1.3\e{-5}}}
&-1.7\e{-7}\std{2.3\e{-7}} 
& 3.777\std{0.103} \\
&(8.5\e{-5}\std{1.0\e{-5}})  
& (7.1\e{-5}\std{1.0\e{-5}}) 
& & (5.3\e{-3}\std{3.9\e{-4}}) 
& (127.105\std{5.372}) \\

\thinrule
\multirow{2}{*}{Projection (qpth)} &  9.7\e{-14}\std{3.5\e{-16}}
& 1.2\e{-14}\std{4.6\e{-16}}
& \multirow{2}{*}{3.7\e{-3}\std{1.7\e{-3}}} 
& 2.2\e{-4}\std{1\e{-4}}  
& 13.71\std{1.07}\\
&(1.4\e{-13}\std{3.4\e{-15}}) 
& (4.1\e{-13}\std{2.2\e{-15}}) 
&&(0.02\std{5.3\e{-3}})
& (225.063\std{43.585}) \\

\thinrule
\rowcolor{darkblue!4}
& 1.3\e{-4}\std{1.1\e{-5}} 
& 1.6\e{-6}\std{4.4\e{-7}} 
& \multirow{2}{*}{0.015\std{0.002}}
& 5.7\e{-3}\std{9.7\e{-4}} 
& 0.121\std{0.038} \\
\rowcolor{darkblue!4}
\multirow{-2}{*}{FSNet} & (7.7\e{-4}\std{1.9\e{-4}}) 
& (7.2\e{-5}\std{3.2\e{-5}}) 
& \multirow{-2}{*}{0.015\std{0.002}}  
& (0.038\std{7.7\e{-4}})  
& (180.389\std{22.505}) \\

\thinrule
\multirow{2}{*}{Penalty} & 0.219\std{5.8\e{-3}} 
& 0.023\std{3.3\e{-3}}
& \multirow{2}{*}{-0.029\std{0.013}} 
& -0.074\std{0.013}  
& 0.043\std{2.9\e{-3}}  \\
&(5.4\e{-5}\std{2.1\e{-5}}) 
&(0.227\std{0.039}) 
& & (0.028\std{9.8\e{-3}}) 
& (1.594\std{0.051}) \\

\thinrule
\multirow{2}{*}{Adaptive Penalty} & 0.221\std{7.6\e{-3}}   
&0.022\std{1.6\e{-3}} 
& \multirow{2}{*}{0.974\std{0.059}} 
& 0.782\std{0.064}  
& 0.042\std{1.7\e{-3}} \\ 
&(0.371\std{2.9\e{-3}}) 
& (0.211\std{0.019})
&&(1.213\std{0.074})
&(1.549\std{0.011}) \\

\thinrule
\multirow{2}{*}{DC3} &  
4.9\e{-13}\std{7.8\e{-16}}
& 0.133\std{0.011}
& \multirow{2}{*}{-0.035\std{0.016}} 
& -0.087\std{0.019}  
& 0.058\std{3.5\e{-3}}\\
&(5.8\e{-13}\std{6.8\e{-15}}) 
& (0.259\std{0.021}) 
&&(0.029\std{0.021})
& (53.151\std{0.558}) \\

%%%%%%%%%%%%%%%%%%%%%%%%%%%%%%%%%%%%%%%%%%%%%%
\midrule
\multicolumn{6}{c}{\textbf{Convex QCQP: \(n=100, n_{\mathrm{eq}}=50, n_{\mathrm{ineq}}=50\)}} \\
\midrule
\multirow{2}{*}{Solver} & 8.4\e{-7}\std{1.9\e{-8}} 
& 1.5\e{-3}\std{2.0\e{-5}} 
& \multirow{2}{*}{--}
& \multirow{2}{*}{--} 
& 46.619\std{0.577}  \\
& (2.4\e{-5}\std{5.8\e{-6}}) 
& (7.7\e{-3}\std{ 0.00}) 
&&  
& (1118.149\std{3.535}) \\


\thinrule
\rowcolor{darkblue!4}
 & 1.2\e{-4}\std{1.2\e{-5}}
& 1.0\e{-6}\std{2.2\e{-7}} 
& 
& 7.8\e{-3}\std{1.1\e{-3}} 
& 0.448\std{0.016} \\
\rowcolor{darkblue!4}
\multirow{-2}{*}{FSNet}
& (1.5\e{-3}\std{5.7\e{-4}}) 
& (6.4\e{-5}\std{2.4\e{-5}}) 
& \multirow{-2}{*}{0.035\std{3.2\e{-3}}} 
&(0.359\std{0.069})  
& (182.156\std{12.827}) \\

\thinrule
\multirow{2}{*}{Penalty} 
& 0.191\std{5.9\e{-3}}
& 0.031\std{0.01}
& \multirow{2}{*}{0.038\std{0.062}} 
& -0.237\std{0.022}  
& 0.048\std{3.9\e{-3}}  \\

&(0.355\std{0.014}) 
&(0.521\std{0.019}) 
& 
& (0.397\std{0.2}) 
& (1.707\std{6.8\e{-3}}) \\

\thinrule
\multirow{2}{*}{Adaptive Penalty} 
& 0.179\std{7.6\std{-3}}
&0.028\std{2.1\e{-3}} 
& \multirow{2}{*}{0.274\std{0.042}} 
& -0.119\std{0.067}  
& 0.047\std{2.8\e{-3}} \\

&(0.333\std{7.6\e{-3}}) 
& (0.497\std{0.075})
&
&(0.797\std{0.079})
&(1.703\std{0.02}) \\

\thinrule
\multirow{2}{*}{DC3} 
& 3.0\e{-13}\std{4.9\e{-16}}
& 0.120\std{0.032}
& \multirow{2}{*}{-0.018\std{0.037}} 
& -0.198\std{0.041}  
&  0.163\std{2.3\e{-3}}\\

&(3.8\e{-13}\std{6.4\e{-15}}) 
&(0.602\std{0.072}) 
&
&(0.259\std{0.039})
&(79.234\std{6.847}) \\

%%%%%%%%%%%%%%%%%%%%%%%%%%%%%%%%%%%%%%%
\midrule
\multicolumn{6}{c}{\textbf{Convex SOCP: \(n=100, n_{\mathrm{eq}}=50, n_{\mathrm{ineq}}=50\)}} 
\\
\midrule
\multirow{2}{*}{Solver} & 3.1\e{-14}\std{4.2\e{-17}} 
& 1.3\e-11 \std{1.2\e{-11}} & \multirow{2}{*}{--}
& \multirow{2}{*}{--} & 332.109\std{2.186} \\
& (2.2\e{-13}\std{0.00}) 
& (1.0\e{-8}\std{8.8\e{-9}}) 
&&  
& (3702.308\std{9.56}) \\

\thinrule
\multirow{2}{*}{Reduced Solver} & 5.9\e{-7}\std{1.3\e{-8}} 
& 1.8\e{-4} \std{4.1\e{-6}} 
& \multirow{2}{*}{4.8\e{-5}\std{2.1\e{-6}}}
& -1.2\e{-3}\std{4.5\e{-5}} & 330.143\std{1.687} \\
&( 4.8\e{-6}\std{3.3\e{-7}})  
& (2.0\e{-3}\std{8.7\e{-5}}) 
& & (1.4\e{-3}\std{1.1\e{-4}}) 
& (1895.989\std{4.592}) \\

\thinrule
\rowcolor{darkblue!4}
& 1.3\e{-4}\std{1.2\e{-5}} 
& 1.7\e{-8}\std{1.3}\e{-8}
& 
& 0.046\std{3.6\e{-3}} 
&  0.295\std{8.0\e{-3}}\\
\rowcolor{darkblue!4}
\multirow{-2}{*}{FSNet} 
& (5.5\e{-4}\std{9.9\e{-5}}) 
& (4.9\e{-6}\std{1.5\e{-6}}) 
& \multirow{-2}{*}{0.159\std{5.0\e{-3}}}
&(1.392\std{0.376}) 
& (208.399\std{9.935}) \\

\thinrule
\multirow{2}{*}{Penalty} 
& 0.098\std{5.2\e{-3}} 
& 0.015\std{3.2\e{-3}}
& \multirow{2}{*}{0.024\std{0.018}} 
& -0.235\std{0.107}  
& 0.049\std{3.4\e{-3}}\\
&(0.139\std{3.6\e{-3}}) 
&(0.309\std{0.053}) 
& 
& (0.338\std{0.108}) 
& (1.919\std{0.236}) \\

\thinrule
\multirow{2}{*}{Adaptive Penalty} 
&0.064\std{2.5\e{-3}}
&0.011\std{1.3\e{-3}}
& \multirow{2}{*}{0.256\std{0.126}} 
& -0.171\std{0.146}  
& 0.047\std{1.4\e{-3}} \\
&(0.136\std{7.2\e{-3}}) 
&(0.305\std{0.147})
&
&(0.808\std{0.359})
&(1.891\std{2.239}) \\

\thinrule
\multirow{2}{*}{DC3} 
&7.8\e{-14}\std{6.9\e{-17}}
&0.029\std{6.0\e{-3}}
& \multirow{2}{*}{0.053\std{0.011}}
& -0.103\std{9.5\e{-3}} 
& 0.114\std{5.4\e{-3}} \\
&(1.1\e{-13}\std{9.8\e{-16}}) 
&(0.325\std{0.049}) 
&
&(0.304\std{0.064})
&(78.83\std{5.254}) \\

\bottomrule
\end{tabular}
\caption{Test results of on 2000 instances of smooth convex problems across 3 random seeds.}
\label{tab: convex results}
\end{table*}

\begin{figure}[h]
    \centering
    \includegraphics[width=0.7\linewidth]{Figures/batch_times.pdf}
    \caption{Computational times of FSNet on 2000 instances of smooth convex and nonconvex (NC) problems with varying batch sizes.}
    \label{fig: time wrt batch}
\end{figure}

\begin{figure}[h]
    \centering
    \includegraphics[width=0.7\linewidth]{Figures/cpu_usage.pdf}
    \caption{CPU utilization over time for parallel solver runs on smooth convex and nonconvex problems.}
    \label{fig: CPU usage}
\end{figure}

\newpage
\subsection{Smooth nonconvex problems}
Figure~\ref{fig: nonconvex results} and Table~\ref{tab: nonconvex results} present the detailed results for nonconvex problems. 
Figure~\ref{fig: time wrt batch} and \ref{fig: CPU usage} also show the computational time of FSNet with different batch sizes and CPU utilization on nonconvex cases, respectively.
The patterns are the same as those observed in the smooth convex problems. 

\begin{figure}[h]
    \centering
    \includegraphics[width=0.95\linewidth]{Figures/nonconvex_results.pdf}
    \caption{Test results on 2000 instances of smooth nonconvex problems. FSNet consistently obtains near-zero constraint violations and small optimality gaps across the problem classes, whereas the baseline methods incur significant constraint violations.}
    \label{fig: nonconvex results}
\end{figure}

\begin{table*}[t]
\centering
\scriptsize
\setlength{\tabcolsep}{5pt}
\begin{tabular}{c c c cc c}
\toprule
\multirow{2}{*}{\textbf{Method}}
& \multicolumn{1}{c}{\textbf{Equality Vio.}}
& \multicolumn{1}{c}{\textbf{Inequality Vio.}} 
& \multicolumn{2}{c}{\textbf{Optimality Gap ($\%$)}} 
& \multicolumn{1}{c}{\textbf{Runtime (s)}} \\
& Mean (Max) 
& Mean (Max)
& Mean  &Min (Max)
& Batch (sequential)\\

\midrule
\multicolumn{6}{c}{\textbf{Nonconvex QP: \(n=100, n_{\mathrm{eq}}=100, n_{\mathrm{ineq}}=100\)}} \\
\midrule
\multirow{2}{*}{Solver} & 5.5\e{-14}\std{1.1\e{-16}} 
& 1.8\e{-9}\std{7.7\e{-11}} 
& \multirow{2}{*}{--}
& \multirow{2}{*}{--} 
& 12.699\std{0.692} \\
& (8.3\e{-14}\std{5.1\e{-15}}) 
& (1.6\e{-8}\std{8.6\e{-10}}) 
&&  
& (470.846\std{0.991}) \\
\thinrule
\rowcolor{darkblue!4}
& 9.3\e{-5}\std{1.9\e{-6}} 
& 7.2\e{-7}\std{8.7\e{-8}} 
&&-0.341\std{0.24} 
& 0.146\std{0.038} \\
\rowcolor{darkblue!4}
\multirow{-2}{*}{FSNet} & (1.4\e{-3}\std{2\e{-4}}) 
&(4.4\e{-5}\std{1.1\e{-5}}) 
&\multirow{-2}{*}{0.146\std{0.012}}  
&(8.885\std{0.592})  
& (171.226\std{10.671}) \\

\midrule
\multicolumn{6}{c}{\textbf{Nonconvex QCQP: \(n=100, n_{\mathrm{eq}}=50, n_{\mathrm{ineq}}=50\)}} \\
\midrule
\multirow{2}{*}{Solver} & 6.8\e{-14}\std{1.6\e{-16}} 
& 8.9\e{-9}\std{2.1\e{-10}}
& \multirow{2}{*}{--}
& \multirow{2}{*}{--} 
& 216.772\std{2.650}  \\
& (9.9\e{-14}\std{3.5\e{-15}}) 
& (5.4\e{-8}\std{1.6\e{-9}})
&& 
&(11169.226\std{102.22}) \\


\thinrule
\rowcolor{darkblue!4}
 & 9.7\e{-5}\std{3.9\e{-6}}
& 2.9\e{-6}\std{2.4\e{-7}} 
& 
& -19.436\std{1.767} 
& 0.221\std{0.028} \\
\rowcolor{darkblue!4}
\multirow{-2}{*}{FSNet}
& (6.5\e{-6}\std{2.9\e{-6}}) 
& (1.1\e{-4}\std{1.4\e{-5}}) 
& \multirow{-2}{*}{-0.668\std{0.650}}
&(8.809\std{0.942})  
& (303.139\std{25.742}) \\

%%%%%%%%%%%%%%%%%%%%%%%%%%%%%%%%%%%%%%%
\midrule
\multicolumn{6}{c}{\textbf{Nonconvex SOCP: \(n=100, n_{\mathrm{eq}}=50, n_{\mathrm{ineq}}=50\)}} 
\\
\midrule
\multirow{2}{*}{Solver} & 8.3\e{-14}\std{1.3\e{-16}} 
& 9.9\e{-9}\std{3.0\e{-10}} 
& \multirow{2}{*}{--}
& \multirow{2}{*}{--} 
& 1216.105\std{7.923} \\
& (1.3\e{-13}\std{3.3\e{-15}}) 
& (7.6\e{-8}\std{7.6\e{-10}})
&&
&(33765.581\std{170.102}) \\

\thinrule
\rowcolor{darkblue!4}
& 5.8\e{-5}\std{3.5\e{-5}} 
& 5.2\e{-7}\std{1.5\e{-7}}
& 
& -1.3\e{6}\std{9.0\e{5}}
& 0.406\std{0.051}  \\
\rowcolor{darkblue!4}
\multirow{-2}{*}{FSNet} 
& (3.4\e{-3}\std{1.4\e{-4}}) 
& (2.5\e{-4}\std{1.6\e{-3}}) 
& \multirow{-2}{*}{-2.3\e{3}\std{0.7\e{3}}}
&(67.138\std{3.265})
& (288.161\std{9.598}) \\
\bottomrule
\end{tabular}
\caption{Test results of FSNet on 2000 instances of smooth nonconvex problems.
}
\label{tab: nonconvex results}
\end{table*}


\newpage
\subsection{Nonsmooth nonconvex problems}

Table~\ref{tab: nonsmooth nonconvex results} shows the detailed results of nonsmooth nonconvex problems (see Appendix~\ref{appendix: problem formulations} for formulations). 
As shown in this table, FSNet achieves near-zero constraint violations across nonsmooth nonconvex problems and achieves 152-3316$\times$ speedups in the batch setting and 7-80$\times$ speedups in the sequential setting compared to the traditional solver, which is also visualized in Figure~\ref{fig: NSNC times}.

Notably, FSNet can attain lower‐cost local solutions than the solver, as evidenced by the negative optimality gaps in Table~\ref{tab: nonsmooth nonconvex results} and by the distributions of objective values in Figure~\ref{fig: nonsmooth nonconvex results}.

\begin{table*}[h]
\centering
\scriptsize
\setlength{\tabcolsep}{5pt}
\begin{tabular}{c c c cc cc@{}}
\toprule
\multirow{2}{*}{\textbf{Method}}
& \multicolumn{1}{c}{\textbf{Equality Vio.}}
& \multicolumn{1}{c}{\textbf{Inequality Vio.}} 
& \multicolumn{2}{c}{\textbf{Optimality Gap ($\%$)}} 
& \multicolumn{1}{c}{\textbf{Runtime (s)}} \\

& Mean (Max)
& Mean (Max)
& Mean  & Min (Max)
& Batch (sequential)\\

\midrule
\multicolumn{6}{c}{\textbf{Nonsmooth Nonconvex QP: \(n=100, n_{\mathrm{eq}}=100, n_{\mathrm{ineq}}=100\)}} \\
\midrule
\multirow{2}{*}{Solver} & 5.4\e{-14}\std{8.2\e{-17}} 
& 5.3\e{-10}\std{5.6\e{-12}} 
& \multirow{2}{*}{--}
& \multirow{2}{*}{--} 
& 19.161\std{0.566} \\
& (7.9\e{-14}\std{6.9\e{-16}}) 
& (1.2\e{-8}\std{2.4\e{-10}}) 
&&  
& (1120.035\std{8.607}) \\
\thinrule
\rowcolor{darkblue!4}
& 9.0\e{-5}\std{3.8\e{-6}} 
& 6.3\e{-7}\std{3.2\e{-8}} 
& 
&-1.081\std{0.275} 
& 0.126\std{0.024} \\
\rowcolor{darkblue!4}
\multirow{-2}{*}{FSNet} & (1.5\e{-3}\std{1.7\e{-4}}) 
& (4.2\e{-5}\std{6.7\e{-6}}) 
& \multirow{-2}{*}{0.109\std{8.5\e{-3}}} 
& (8.128\std{0.593})  
& (156.553\std{4.761}) \\

%%%%%%%%%%%%%%%%%%%%%%%%%%%%%%%%%%%%%%%%%%%%%%
\midrule
\multicolumn{6}{c}{\textbf{Nonconvex QCQP: \(n=100, n_{\mathrm{eq}}=50, n_{\mathrm{ineq}}=50\)}} \\
\midrule
\multirow{2}{*}{Solver} & 6.8\e{-14}\std{2.4\e{-16}} 
& 4.9\e{-9}\std{3.7\e{-11}}
& \multirow{2}{*}{--}
& \multirow{2}{*}{--} 
& 288.768\std{0.329}  \\
& (1.0\e{-13}\std{1.4\e{-15}}) 
& (4.5\e{-8}\std{3.5\e{-9}}) 
&&  
& (15092.158\std{6116.515}) \\

\thinrule
\rowcolor{darkblue!4}
& 9.1\e{-5}\std{4.8\e{-6}}
& 2.9\e{-6}\std{1.69\e{-7}} 
& 
& -27.299\std{2.573} 
& 0.539\std{0.018} \\
\rowcolor{darkblue!4}
\multirow{-2}{*}{FSNet} 
& (6.5\e{-6}\std{2.9\e{-6}}) 
& (1.1\e{-4}\std{1.4\e{-5}}) 
& \multirow{-2}{*}{-3.437\std{0.661}} 
& (9.466\std{0.919})  
& (247.819\std{37.821}) \\

%%%%%%%%%%%%%%%%%%%%%%%%%%%%%%%%%%%%%%%
\midrule
\multicolumn{6}{c}{\textbf{Nonsmooth Nonconvex SOCP: \(n=100, n_{\mathrm{eq}}=50, n_{\mathrm{ineq}}=50\)}} 
\\
\midrule
\multirow{2}{*}{Solver} & 7.8\e{-14}\std{1.3\e{-16}} 
& 3.4\e{-9}\std{1.5\e{-10}} 
& \multirow{2}{*}{--}
& \multirow{2}{*}{--} 
& 1363.465\std{5.109} \\
& (1.2\e{-13}\std{2.5\e{-15}}) 
& 5.9\e{-8}\std{2.7\e{-9}} 
&&  
& (21642.764\std{81.569}) \\

\thinrule
\rowcolor{darkblue!4}
& 6.2\e{-5}\std{9.4\e{-6}} 
& 4.2\e{-7}\std{1.2\e{-7}}
& 
& -9.6\e{5}\std{7.6\e{5}}
& 0.411\std{0.019}  \\
\rowcolor{darkblue!4}
\multirow{-2}{*}{FSNet} 
& (2.5\e{-3}\std{5.2\e{-4}}) 
& (6.1\e{-5}\std{6.0\e{-6}}) 
& \multirow{-2}{*}{-1.2\e{3}\std{9.43\e{2}}}
&(658.5\std{831.8})
& (267.389\std{6.987}) \\
\bottomrule
\end{tabular}
\caption{Test results of FSNet on 2000 instances of nonsmooth nonconvex problems.}
\label{tab: nonsmooth nonconvex results}
\end{table*}

\begin{figure}[h]
    \centering
    \includegraphics[width=0.8\linewidth]{Figures/histograms_nonsmooth_nonconvex.pdf}
    \caption{Distribution of the objective values in nonsmooth nonconvex problems.}
    \label{fig: nonsmooth nonconvex results}
\end{figure}

\begin{figure}[h]
    \centering
    \includegraphics[width=0.7\linewidth]{Figures/NSNC_times.pdf}
    \caption{Computational time on 2000 instances of nonsmooth nonconvex (NSNC) problems.}
    \label{fig: NSNC times}
\end{figure}

\newpage
\subsection{FSNet with truncated unrolled differentiation}
This section evaluates the impact of the truncation depth on the performance of FSNet (see Section \ref{sec: improvement}). We set the maximum number of feasibility‐seeking iterations to 50, and vary the truncation depth from 0 to 50. The results on the smooth convex problems are shown in Table~\ref{tab: different gradient-tracked iterations} and Figure~\ref{fig: time_comparison_diff}.
\begin{figure}[h]
    \centering
    \includegraphics[width=0.95\linewidth]{Figures/time_comparison_diff.pdf}
    \caption{Training time and batch inference time of FSNet with different truncation depths.}
    \label{fig: time_comparison_diff}
\end{figure}

When $K'=0$, no iterations of the feasibility-seeking procedure are included in the computational graph, so it acts as a non-differentiable pass-through layer. 
This yields the smallest graph and therefore the shortest training time, but incurs a large optimality gap, indicating the network cannot learn the true solutions of Problem~\eqref{prob: nonlinear problem}. 
As $K'$ increases, the training is successful with small optimality gaps (and near-zero constraint violations), yet training time grows due to the larger graph and additional backpropagation cost (Figure~\ref{fig: time_comparison_diff}).
Importantly, for $K'$ between 10 and 50, both total constraint violation and optimality gap remain essentially unchanged. 
This means that a modest value of $K'$ is sufficient to render the bias caused by the gradient truncation negligible, which agrees with our analysis in Appendix~\ref{subsec: truncated diff analysis}. 
This suggests a practical tip that, rather than fully gradient-tracking all iterations of the feasibility-seeking procedure, we only need to track a few of them while maintaining high performance for FSNet and reducing the training time. 

Since truncation affects only the backward pass during training, inference times are virtually identical for all choices of $K' \neq 0$.


\begin{table*}[h]
\centering
\scriptsize
\setlength{\tabcolsep}{5pt}
\begin{tabular}{cccccc}
\toprule
\multirow{2}{*}{\textbf{Tracked Iter. $K'$}}
& \multicolumn{1}{c}{\textbf{Total vio.}}
& \multicolumn{1}{c}{\textbf{Optimality Gap ($\%$)}} 
& \multicolumn{1}{c}{\textbf{Runtime (s)}} 
& \multirow{2}{*}{\textbf{Training time (h)}} \\

& Mean (Max) 
& Mean (Max) 
& Batch (Sequential)\\
\midrule
\multicolumn{5}{c}{\textbf{Convex QP: \(n=100, n_{\mathrm{eq}}=100, n_{\mathrm{ineq}}=100\)}} \\
\midrule
\multirow{2}{*}{0} 
& 4.4\e{-5}\std{1\e{-6}} 
& 14.143\std{0.024} 
& 0.174\std{0.046}
& \multirow{2}{*}{0.348\std{0.009}} \\
& (4.2\e{-4}\std{4.2\e{-5}}) 
& (16.314\std{0.148}) 
& (296.889\std{64.220}) &  \\
\thinrule
\multirow{2}{*}{10} & 7.5\e{-5}\std{6.2\e{-6}} 
& 0.015\std{0.001} 
& 0.134\std{0.036}
& \multirow{2}{*}{0.490\std{0.015}} \\
&(6.1\e{-4}\std{1\e{-4}}) 
& (0.050\std{0.005}) 
& (168.696\std{7.896}) &  \\

\thinrule
\multirow{2}{*}{20}
&6.8\e{-5}\std{8.3\e{-6}} 
& 0.011\std{0.001} 
& 0.149\std{0.039} 
& \multirow{2}{*}{0.668\std{0.037}}\\ 
&(5.5\e{-4}\std{4.6\e{-5}}) 
& (0.034\std{0.002}) 
& (162.488\std{37.734}) & \\ 

\thinrule
\multirow{2}{*}{30}
& 6.6\e{-5}\std{5.5\e{-6}} 
& 0.015\std{0.002} 
& 0.124\std{0.033} 
& \multirow{2}{*}{0.695\std{0.008}}
\\ 
& (9.2\e{-4}\std{1.9\e{-4}}) 
& (0.046\std{0.008}) 
& (166.242\std{45.394}) & \\ 

\thinrule
\multirow{2}{*}{40}
& 6.6\e{-5}\std{5.5\e{-6}} 
& 0.015\std{0.002} 
& 0.125\std{0.040} 
& \multirow{2}{*}{0.682\std{0.011}} \\ 
& (9.2\e{-4}\std{1.9\e{-4}}) 
& (0.046\std{0.008}) 
& (136.559\std{3.089}) & \\

\thinrule
\multirow{2}{*}{50}
& 6.6\e{-5}\std{5.5\e{-6}} 
& 0.015\std{0.002} 
& 0.125\std{0.033} 
& \multirow{2}{*}{0.691\std{0.019}}\\ 
& (9.2\e{-4}\std{1.9\e{-4}}) 
& (0.046\std{0.008}) 
& (176.246\std{49.894}) & \\

%%%%%%%%%%%%%%%%%%%%%%%%%%%%%%%%%%%%%%%%%%%%%%
\midrule
\multicolumn{5}{c}{\textbf{Convex QCQP: \(n=100, n_{\mathrm{eq}}=50, n_{\mathrm{ineq}}=50\)}} \\
\midrule
\multirow{2}{*}{0} 
& 7.2\e{-5}\std{4.5\e{-7}} 
& 37.174\std{0.148} 
& 0.648\std{0.003}
& \multirow{2}{*}{0.461\std{0.011}}
\\
&(7.9\e{-4}\std{9.7\e{-5}}) 
& (45.584\std{0.395}) 
& (284.443\std{35.772}) &  \\
\thinrule
\multirow{2}{*}{10} 
& 6.2\e{-5}\std{2.8\e{-6}} 
& 0.032\std{0.002} 
& 0.475\std{0.008}
& \multirow{2}{*}{0.536\std{0.037}}
 \\
&(8.1\e{-4}\std{6.6\e{-5}}) 
& (0.405\std{0.062}) 
& (193.722\std{28.776}) &  \\
\thinrule
\multirow{2}{*}{20}
&6.2\e{-5}\std{8.0\e{-6}} 
& 0.032\std{0.002} 
& 0.451\std{0.002} 
& \multirow{2}{*}{0.771\std{0.047}}
\\ 
& (1.1e-03\std{1.8\e{-4}}) 
& (0.411\std{0.052}) 
& (223.545\std{27.455}) & \\ 
\thinrule
\multirow{2}{*}{30}
&6.2\e{-5}\std{5.8\e{-6}} 
& 0.035\std{0.003} 
& 0.457\std{0.009} 
& \multirow{2}{*}{0.757\std{0.028}}
\\ 
& (2.1e-03\std{5.7\e{-4}}) 
& (0.432\std{0.070}) 
& (213.981\std{18.431}) & \\ 
\thinrule
\multirow{2}{*}{40}
& 6.2\e{-5}\std{5.8\e{-6}} 
& 0.035\std{0.003} 
& 0.440\std{0.006} 
& \multirow{2}{*}{0.729\std{0.015}}
\\ 
& (2.1e-03\std{5.7\e{-4}}) 
& (0.432\std{0.070}) 
& (200.018\std{28.512}) & \\
\thinrule
\multirow{2}{*}{50}
& 6.2\e{-5}\std{5.8\e{-6}} 
& 0.035\std{0.003} 
& 0.446\std{0.003} 
& \multirow{2}{*}{0.765\std{0.026}}
\\ 
& (2.1e-03\std{5.7\e{-4}}) 
& (0.432\std{0.070}) 
& (174.800\std{28.764}) & \\ 

%%%%%%%%%%%%%%%%%%%%%%%%%%%%%%%%%%%%%%%
\midrule
\multicolumn{5}{c}{\textbf{Convex SOCP: \(n=100, n_{\mathrm{eq}}=50, n_{\mathrm{ineq}}=50\)}} 
\\
\midrule
\multirow{2}{*}{0} 
& 4.8\e{-5}\std{2\e{-5}} 
& 181.38\std{0.092} 
& 0.416\std{0.016}
& \multirow{2}{*}{0.406\std{0.032}}
 \\
& (4.7\e{-4}\std{7.4\e{-5}}) 
& (194.957\std{0.896}) 
& (354.401\std{86.816}) &  \\
\thinrule
\multirow{2}{*}{10}
& 6.6\e{-5}\std{7.3\e{-6}} 
& 0.136\std{0.010} 
& 0.288\std{0.008} 
& \multirow{2}{*}{0.532\std{0.050}}
\\ 
&(6.2\e{-4}\std{6.3\e{-5}}) 
& (1.006\std{0.143}) 
& (252.299\std{56.041}) & \\ 

\thinrule
\multirow{2}{*}{20}
&6.2\e{-5}\std{2.2\e{-6}} 
& 0.186\std{0.020} 
& 0.289\std{0.003} 
& \multirow{2}{*}{0.717\std{0.037}}
\\ 
&(5.6\e{-4}\std{2.2\e{-5}}) 
& (1.914\std{0.439}) 
& (187.842\std{26.580}) & \\ 

\thinrule
\multirow{2}{*}{30}
&6.3\e{-5}\std{6.1\e{-6}} 
& 0.159\std{0.005} 
& 0.336\std{0.036} 
& \multirow{2}{*}{0.780\std{0.024}}
\\ 
& (6.7\e{-4}\std{9.9\e{-5}}) 
& (1.889\std{0.376}) 
& (272.440\std{52.514}) & \\

\thinrule
\multirow{2}{*}{40}
& 6.3\e{-5}\std{6.1\e{-6}} 
& 0.159\std{0.005} 
& 0.297\std{0.010} 
& \multirow{2}{*}{0.804\std{0.015}}
\\ 
& (6.7\e{-4}\std{9.9\e{-5}}) 
& (1.889\std{0.376}) 
& (214.608\std{17.089}) & \\

\thinrule
\multirow{2}{*}{50}
&6.3\e{-5}\std{6.1\e{-6}} 
& 0.159\std{0.005} 
& 0.297\std{0.008} 
& \multirow{2}{*}{0.771\std{0.043}}\\ 
& (6.7\e{-4}\std{9.9\e{-5}}) 
& (1.889\std{0.376}) 
& (259.888\std{30.394}) & \\

\bottomrule
\end{tabular}
\caption{Numerical results on 2000 instances of smooth convex problems with varying numbers of tracked iterations.}
\label{tab: different gradient-tracked iterations}
\end{table*}

\newpage
\subsection{FSNet with different values of $\rho$}
This section investigates the impact of the value of penalty weight $\rho$ on the performance of FSNet (see Section~\ref{sec: fs}).
Table~\ref{tab: different rho} shows that the constraint violations are near-zero for all choices of $\rho$, since the feasibility-seeking step always tries to enforce the solution's feasibility regardless of the NN's prediction.
This is an advantage of FSNet, as the feasibility is guaranteed by a well-designed feasibility-seeking step, rather than being left to the NN alone.
While the optimality gaps are similar across different $\rho$, both training and inference times decrease as $\rho$ increases (Figure~\ref{fig: time_comparison_rho}).
Actually, a larger penalty weight strongly penalizes the distance between the solutions before and after the feasibility-seeking step, driving the NN to produce outputs that are closer to the feasible set. 
This reduces the number of iterations required by the feasibility-seeking step, leading to lower computational time for the entire process. 
For example, setting $\rho=50$ makes the training 1.4-1.6$\times$ faster and inference 1.3-1.7$\times$ faster compared to $\rho=0$. 
However, excessively large values of $\rho$ may cause the optimization to overemphasize feasibility at the expense of optimality, potentially slowing convergence of the network parameters.

\begin{table*}[h]
\centering
\scriptsize
\setlength{\tabcolsep}{5pt}
\begin{tabular}{cccccc}
\toprule
\multirow{2}{*}{\textbf{$\rho$}}
& \multicolumn{1}{c}{\textbf{Total vio.}}
& \multicolumn{1}{c}{\textbf{Optimality Gap ($\%$)}} 
& \multicolumn{1}{c}{\textbf{Runtime (s)}} 
& \multicolumn{1}{c}{\textbf{Training time (h)}} \\

& Mean (Max) 
& Mean (Max) 
& Batch (Sequential)\\
\midrule
\multicolumn{5}{c}{\textbf{Convex QP: \(n=100, n_{\mathrm{eq}}=100, n_{\mathrm{ineq}}=100\)}} \\
\midrule
\multirow{2}{*}{0} 
&5.8e-05\std{6.2e-06} 
& 0.008\std{0.001} 
& 0.183\std{0.055} 
& \multirow{2}{*}{1.075\std{0.027}}
\\ 
&(5.8e-04\std{1.0e-04}) 
& (0.029\std{0.003}) 
& (261.704\std{13.474}) & \\ 

\thinrule
\multirow{2}{*}{0.5} 
&6.5e-05\std{1.0e-05} 
& 0.009\std{0.001} 
& 0.150\std{0.034} 
& \multirow{2}{*}{0.827\std{0.076}}
\\ 
&(9.5e-04\std{2.4e-04}) 
& (0.031\std{0.004}) 
& (236.686\std{50.218}) & \\

\thinrule
\multirow{2}{*}{2.5} 
&6.6e-05\std{8.6e-06} 
& 0.011\std{0.002} 
& 0.183\std{0.014} 
& \multirow{2}{*}{0.747\std{0.044}}\\ 
&(7.8e-04\std{1.4e-04}) 
& (0.041\std{0.004}) 
& (196.120\std{47.639}) & \\

\thinrule
\multirow{2}{*}{5.0}
&6.6e-05\std{5.5e-06} 
& 0.015\std{0.002} 
& 0.142\std{0.027} 
& \multirow{2}{*}{0.687\std{0.020}}
\\ 
&(9.2e-04\std{1.9e-04}) 
& (0.046\std{0.008}) 
& (186.968\std{34.144}) & \\ 
\thinrule
\multirow{2}{*}{10.0}
&7.1e-05\std{6.9e-06} 
& 0.014\std{0.001} 
& 0.117\std{0.029} 
& \multirow{2}{*}{0.676\std{0.022}}
\\ 
&(9.6e-04\std{2.0e-04}) 
& (0.042\std{0.004}) 
& (152.014\std{22.923}) & \\

\thinrule
\multirow{2}{*}{20.0}
&6.7e-05\std{2.4e-06} 
& 0.011\std{0.002} 
& 0.133\std{0.050} 
& \multirow{2}{*}{0.649\std{0.013}}
\\ 
&(1.2e-03\std{2.5e-04}) 
& (0.042\std{0.003}) 
& (165.496\std{23.169}) & \\
\thinrule
\multirow{2}{*}{50.0}
&5.9e-05\std{1.1e-05} 
& 0.018\std{0.006} 
& 0.104\std{0.013} 
& \multirow{2}{*}{0.661\std{0.012}}
\\ 
&(9.3e-04\std{2.5e-04}) 
& (0.074\std{0.016}) 
& (156.569\std{53.576}) & \\ 

%%%%%%%%%%%%%%%%%%%%%%%%%%%%%%%%%%%%%%%%%%%%%%
\midrule
\multicolumn{5}{c}{\textbf{Convex QCQP: \(n=100, n_{\mathrm{eq}}=50, n_{\mathrm{ineq}}=50\)}} \\
\midrule
\multirow{2}{*}{0} 
&7.2e-05\std{2.7e-06} 
& 0.023\std{0.002} 
& 0.576\std{0.004} 
& \multirow{2}{*}{1.039\std{0.037}}
\\ 
&(7.3e-04\std{1.6e-04}) 
& (0.315\std{0.064}) 
& (257.309\std{15.105}) & \\
\thinrule
\multirow{2}{*}{0.5} 
&5.7e-05\std{1.0e-05} 
& 0.028\std{0.002} 
& 0.487\std{0.015} 
& \multirow{2}{*}{0.811\std{0.010}}
\\ 
&(1.0e-03\std{1.5e-04}) 
& (0.412\std{0.075}) 
& (216.307\std{32.152}) & \\
\thinrule
\multirow{2}{*}{2.5} 
&7.1e-05\std{7.8e-06} 
& 0.030\std{0.001} 
& 0.464\std{0.004} 
& \multirow{2}{*}{0.771\std{0.024}}
\\ 
&(1.6e-03\std{3.6e-04}) 
& (0.411\std{0.073}) 
& (187.807\std{11.011}) & \\
\thinrule
\multirow{2}{*}{5.0}
&6.2e-05\std{5.8e-06} 
& 0.035\std{0.003} 
& 0.451\std{0.008} 
& \multirow{2}{*}{0.743\std{0.014}}
\\ 
&(2.1e-03\std{5.7e-04}) 
& (0.432\std{0.070}) 
& (180.282\std{32.059}) & \\ 
\thinrule
\multirow{2}{*}{10.0}
&6.1e-05\std{6.6e-06} 
& 0.038\std{0.003} 
& 0.450\std{0.014} 
& \multirow{2}{*}{0.715\std{0.006}}
\\ 
&(1.5e-03\std{1.2e-04}) 
& (0.417\std{0.058}) 
& (222.192\std{56.348}) & \\
\thinrule
\multirow{2}{*}{20.0}
&5.8e-05\std{6.8e-06} 
& 0.041\std{0.007} 
& 0.437\std{0.009} 
& \multirow{2}{*}{0.762\std{0.052}}
\\ 
&(2.4e-03\std{6.9e-04}) 
& (0.466\std{0.081}) 
& (240.089\std{33.913}) & \\
\thinrule
\multirow{2}{*}{50.0}
&5.3e-05\std{1.5e-06} 
& 0.049\std{0.016} 
& 0.447\std{0.015} 
& \multirow{2}{*}{0.716\std{0.031}}
\\ 
&(2.6e-03\std{7.5e-04}) 
& (0.486\std{0.088}) 
& (225.130\std{36.618}) & \\

%%%%%%%%%%%%%%%%%%%%%%%%%%%%%%%%%%%%%%%
\midrule
\multicolumn{5}{c}{\textbf{Convex SOCP: \(n=100, n_{\mathrm{eq}}=50, n_{\mathrm{ineq}}=50\)}} 
\\
\midrule
\multirow{2}{*}{0} 
&5.3e-05\std{8.2e-06} 
& 0.356\std{0.097} 
& 0.360\std{0.012} 
& \multirow{2}{*}{0.976\std{0.044}}
\\ 
&(7.1e-04\std{1.2e-04}) 
& (1.490\std{0.278}) 
& (310.059\std{71.106}) & \\ 
\thinrule
\multirow{2}{*}{0.5} 
&7.2e-05\std{8.3e-06} 
& 0.180\std{0.012} 
& 0.325\std{0.017} 
& \multirow{2}{*}{0.921\std{0.014}}
\\ 
&(6.6e-04\std{7.6e-05}) 
& (2.347\std{0.514}) 
& (237.029\std{72.636}) & \\
\thinrule
\multirow{2}{*}{2.5} 
&6.6e-05\std{9.2e-06} 
& 0.174\std{0.019} 
& 0.294\std{0.007} 
& \multirow{2}{*}{0.797\std{0.028}}
\\ 
&(7.0e-04\std{1.1e-04}) 
& (0.978\std{0.109}) 
& (206.250\std{15.332}) & \\ 
\thinrule
\multirow{2}{*}{5.0}
&6.3e-05\std{6.1e-06} 
& 0.159\std{0.005} 
& 0.307\std{0.014} 
& \multirow{2}{*}{0.837\std{0.005}}
\\ 
&(6.7e-04\std{9.9e-05}) 
& (1.889\std{0.376}) 
& (281.764\std{18.464}) & \\
\thinrule
\multirow{2}{*}{10.0}
&7.9e-05\std{7.0e-07} 
& 0.174\std{0.020} 
& 0.291\std{0.011} 
& \multirow{2}{*}{0.725\std{0.011}}
\\ 
&(7.7e-04\std{8.2e-05}) 
& (1.365\std{0.251}) 
& (198.642\std{15.576}) & \\
\thinrule
\multirow{2}{*}{20.0}
&6.1e-05\std{4.1e-06} 
& 0.184\std{0.017} 
& 0.287\std{0.013} 
& \multirow{2}{*}{0.714\std{0.041}}
\\ 
&(9.5e-04\std{1.2e-04}) 
& (0.994\std{0.094}) 
& (197.955\std{33.984}) & \\
\thinrule
\multirow{2}{*}{50.0}
&6.3e-05\std{9.1e-06} 
& 0.190\std{0.016} 
& 0.269\std{0.015} 
& \multirow{2}{*}{0.694\std{0.027}}\\ 
&(7.4e-04\std{9.2e-05}) 
& (1.093\std{0.192}) 
& (172.091\std{41.039}) & \\ 
\bottomrule
\end{tabular}
\caption{Numerical results of FSNet on 2000 instances of smooth convex problems with varying $\rho$.}
\label{tab: different rho}
\end{table*}

\begin{figure}[htb]
    \centering
    \includegraphics[width=0.95\linewidth]{Figures/time_comparison_rho.pdf}
    \caption{Training time and batch inference time of FSNet with different values of \(\rho\).}
    \label{fig: time_comparison_rho}
\end{figure}


\clearpage

\subsection{AC Optimal Power Flow}
% \revise{
We evaluate FSNet in solving the AC optimal power flow (ACOPF) in power systems. We use the formulation and dataset in \cite{klamkin2025pglearn}.
% \todo{report number of training vs. validation vs. test samples?}
In the following experiments, we use 10000 samples for each grid size, split into 8000 for training and 2000 for testing.
The parameters $x$ represent the load demands, and the prediction $y_\theta(x)$ of FSNet is the voltage magnitudes and angles of all power system buses.
The Levenberg-Marquardt method is employed for the feasibility-seeking step.
The results for the IEEE 30-bus and 118-bus systems under three random seeds are presented in Table \ref{tab: acopf results}. FSNet achieves small (numerically zero) equality and inequality constraint violations, and significantly outperforms the Penalty method. 
% satisfies equality and inequality constraints (numerical violations on a similar order to that of the IPOPT solver)
Moreover, the optimality gap remains below 1$\%$ in both test cases. In terms of runtime in the sequential setting, FSNet significantly outperforms IPOPT, achieving 3$\times$ speedups on the IEEE 30-bus system and 11$\times$ on the IEEE 118-bus system. The Penalty method, on the other hand, despite its fast inference, suffers from high constraint violations, particularly in terms of maximum violations on the IEEE 118-bus system. Overall, this result is consistent with our previous findings on convex and nonconvex problems.
% }

\begin{table*}[h]
\centering
\scriptsize
\setlength{\tabcolsep}{5pt}
\begin{tabular}{c c c cc cc@{}}
\toprule
\multirow{2}{*}{\textbf{Method}}
& \multicolumn{1}{c}{\textbf{Equality Vio.}}
& \multicolumn{1}{c}{\textbf{Inequality Vio.}} 
& \multicolumn{2}{c}{\textbf{Optimality Gap ($\%$)}} 
& \multicolumn{1}{c}{\textbf{Runtime (s)}} \\

& Mean (Max)
& Mean (Max)
& Mean  & Min (Max)
& Sequential\\
\midrule
\multicolumn{6}{c}{\textbf{IEEE 30-bus: \(n=60, n_{\mathrm{eq}}=60, n_{\mathrm{ineq}}=248\)}} \\
\midrule
\multirow{2}{*}{IPOPT Solver} & 6.7\e{-7}\std{0.0} 
& 5.4\e{-9}\std{0.0} 
& \multirow{2}{*}{--}
& \multirow{2}{*}{--} 
&  \\
&(5.8\e{-6}\std{0.0}) & (1.1\e{-6}\std{0.0}) &&&\multirow{-2}{*}{51.53\std{0.0}}\\

\thinrule
\rowcolor{darkblue!4}
& 1.8\e{-5}\std{7.5\e{-6}} 
& 1.4\e{-8}\std{4.7\e{-9}}
& 
& 0.07\std{1.7\e{-3}} 
&  \\
\rowcolor{darkblue!4}
\multirow{-2}{*}{FSNet}
& (1.6\e{-3}\std{2.1\e{-4}})  &  (6.7\e{-4}\std{1.7\e{-4}}) & \multirow{-2}{*}{0.81\std{0.21}} &(3.61\std{0.92}) & \multirow{-2}{*}{16.43\std{0.13}}\\

\thinrule
Penalty
& 8.2\e{-4}\std{5.3\e{-5}}  
& 3.7\e{-7}\std{1.2\e{-7}} 
& 
& -0.13\std{0.86} 
&  \\
& (0.013\std{1.5\e{-3}}) & (0.026\std{8.5\e{-3}}) & \multirow{-2}{*}{2.24\std{0.53}} & (7.09\std{0.65})   & \multirow{-2}{*}{0.52\std{0.23}}\\

\midrule
\multicolumn{6}{c}{\textbf{IEEE 118-bus: \(n=236, n_{\mathrm{eq}}=236, n_{\mathrm{ineq}}=1196\)}} \\
\midrule
\multirow{2}{*}{IPOPT Solver} & 1.2\e{-6}\std{0.0} 
& 3.1\e{-8}\std{0.0} 
& \multirow{2}{*}{--}
& \multirow{2}{*}{--} 
&  \\
&(2.8\e{-5}\std{0.0})&(1.1\e{-5}\std{0.0})&&&\multirow{-2}{*}{392.83\std{0.0}} \\

\thinrule
\rowcolor{darkblue!4}
& 7.8\e{-5}\std{1.1\e{-5}} 
& 2.4\e{-6}\std{2.3\e{-7}} 
& & -1.67\std{0.92} 
&  \\
\rowcolor{darkblue!4}
\multirow{-2}{*}{FSNet}
&(0.047\std{0.010}) &(0.043\std{8.5\e{-3}})&\multirow{-2}{*}{0.9\std{0.03}}&(4.91\std{0.09})&\multirow{-2}{*}{33.35 \std{0.61}}\\

\thinrule
\multirow{2}{*}{Penalty}
& 2.5\e{-3}\std{3.2\e{-4}} 
& 3.5\e{-5}\std{9.9\e{-6}} 
& 
& -4.67\std{0.17} 
&  \\
&(0.124\std{0.014}) &(0.289\std{0.041})&\multirow{-2}{*}{0.49\std{0.14}}&(1.51\std{0.25})&\multirow{-2}{*}{0.81\std{0.08}}\\


\bottomrule
\end{tabular}
\caption{Test results of FSNet on 2000 instances of ACOPF problems.}
\label{tab: acopf results}
\end{table*}




\subsection{More experiments for DC3}
In \cite{donti2021dc3}, the authors mention that unrolling more correction steps may improve the feasibility of the prediction. To evaluate this, we ran DC3 with different numbers of correction steps on 2000 test instances of smooth convex SOCP.
As shown in Table~\ref{tab: dc3 steps_sensitivity}, increasing the number of correction steps only yields a marginal reduction in inequality violations (due to the default small step size of $10^{-7}$), while significantly increasing runtime.

Next, we evaluate the performance of DC3 with different correction step sizes. 
We fixed the number of maximum correction steps and changed the stepsize to obtain the results in Table~\ref{tab: dc3 lr_sensitivity}. 
Stepsizes up to $10^{-5}$ achieve limited violation reduction, whereas stepsizes of $10^{-4}$ and $10^{-3}$ substantially mitigate violations. However, overly big stepsize (e.g., $10^{-2}$) causes the correction procedure to diverge. Note that all experiments here are in inference time. During training, a large correction stepsize can destabilize the optimization, as noted in \cite{donti2021dc3}.
This demonstrates the DC3's sensitivity to the correction stepsize, whereas our feasibility-seeking step, implemented via L-BGFS with a backtracking line search, mitigates such sensitivity and delivers robust performance across diverse problems using the same hyperparameters. 


\begin{table*}[t]
\footnotesize
\centering
\begin{tabular}{ccccc}
\toprule
Max Correction Steps & Eq. Violation & Ineq. Violation & Optimality Gap (\%) & Time (s) \\
\midrule
10 & 2.30e-14 & 0.024 & 0.051 & 0.095 \\
50 & 2.30e-14 & 0.024 & 0.051 & 0.474 \\
100 & 2.27e-14 & 0.024 & 0.052 & 0.948 \\
200 & 2.28e-14 & 0.023 & 0.052 & 1.896 \\
500 & 2.28e-14 & 0.022 & 0.054 & 4.738 \\
1000 & 2.34e-14 & 0.020 & 0.057 & 9.526 \\
2000 & 2.28e-14 & 0.016 & 0.061 & 19.046 \\
5000 & 2.27e-14 & 0.010 & 0.070 & 48.024 \\
\bottomrule
\end{tabular}
\caption{Sensitivity of DC3 performance to the maximum number of correction steps on 2000 instances of smooth convex SOCP problem, with fixed correction stepsize of $10^{-7}$.}
\label{tab: dc3 steps_sensitivity}
% \end{table*}
\vspace{2em}
% \begin{table*}[t]
% \footnotesize
% \centering
\begin{tabular}{ccccc}
\toprule
Correction Stepsize & Eq. Violation & Ineq. Violation & Objective Gap (\%) & Time (s) \\
\midrule
1e-07 & 2.30e-14 & 0.024 & 0.051 & 0.338 \\
1e-06 & 2.30e-14 & 0.024 & 0.052 & 0.121 \\
1e-05 & 2.28e-14 & 0.020 & 0.056 & 0.095 \\
1e-04 & 2.28e-14 & 0.004 & 0.081 & 0.096 \\
1e-03 & 2.31e-14 & 0.000 & 0.130 & 0.095 \\
1e-02 & -3.52e+12 & 2.89e+30 & 1e58 & 0.095 \\
\bottomrule
\end{tabular}
\caption{Sensitivity of DC3 performance to the correction stepsize on 2000 instances of smooth convex SOCP problem, with fixed maximum correction steps of 10.}
\label{tab: dc3 lr_sensitivity}
\end{table*}




\clearpage
\section{Hyperparameters} \label{appendix: hyperparameters}
For DC3 and Penalty methods, we adopt the hyperparameters as tuned in \cite{donti2021dc3}. 
For the Adaptive Penalty method, we performed a brief manual sweep to identify settings yielding strong empirical performance. For FSNet, we configured the L‐BFGS solver with a sufficiently large maximum‐iteration budget and memory size to guarantee convergence on all the problem classes. We also evaluated several neural network sizes and learning rates to ensure sufficient expressive capacity and stable optimization for all problem classes.
Table~\ref{tab: hyperparameters} lists all hyperparameter values.

\begin{table*}[ht]
\centering
\footnotesize
\setlength{\tabcolsep}{6pt}
\begin{tabular}{l|c}
\toprule
Hyperparameter & Value \\
\midrule
\hspace{-8pt} \textit{Shared}: & \\
Learning rate & 5\e{-4} \\
Learning rate decay & 0.5 \\
% Learning rate decay steps & 2000 \\
Number of hidden layers & 4 \\
Nonlinearity & SiLU \\
Number of hidden neurons per layer & 1024 \\
Train/validation/test ratio & 0.7/0.1/0.2\\
Minibatch size & 512 \\
Random seeds & 2025, 2027, 2029 \\
\midrule
\hspace{-8pt} \textit{FS}: & \\
Equality violation weight & 10 \\
Inequality violation weight & 10 \\
Max L-BFGS iterations & 50 \\
L-BFGS memory & 30 \\
Epoch & 100 \\
Learning rate decay steps & 2000\\

\midrule
\hspace{-8pt} \textit{Penalty}: & \\
Equality violation weight & 50  \\
Inequality violation weight & 50 \\
Epoch & 1000 \\
Learning rate decay steps & 4000\\

\midrule
\hspace{-8pt} \textit{Adaptive Penalty}: & \\
Initial equality violation weight & 30  \\
Initial inequality violation weight & 30 \\
Max equality violation weight & 500  \\
Max inequality violation weight & 500 \\
Increasing rate & 2 \\
Epoch & 1000 \\
Learning rate decay steps & 4000\\
 
\midrule
\hspace{-8pt} \textit{DC3}: & \\
Equality violation weight & 10  \\
Inequality violation weight & 10 \\
Correction stepsize & 1\e{-7} \\
Max correction stepsize & 10 \\
Correction momentum & 0.5 \\
Epoch & 1000 \\
Learning rate decay steps & 4000

\end{tabular}
\caption{Hyperparameters}
\label{tab: hyperparameters}
\end{table*}
%%%%%%%%%%%%%%%%%%%%%%%%%%%%%%%%%%%%%%%%%%%%%%%%%%%%%%%%%%%%




\clearpage
\section{Analysis and Proofs} \label{appendix: proofs}

\subsection{Proof of Theorem \ref{thm: SGD convergence 1}}
\thmSgdConvergence*

\begin{proof} 
    \textbf{Gradient bound of finite-unrolling loss. } For the first part of the theorem, we follow the analysis of the SGD for a $L$-smooth function. 
    By the quadratic upper bound of the $L$-smooth function, we have 
    \begin{align*}
        L(\theta_{t+1})  \leq L(\theta_t) + \langle \nabla L(\theta_t), \theta_{t+1} - \theta_t \rangle  + \frac{L_N}{2} \| \theta_{t+1} - \theta_t \|^2_2.
    \end{align*}
    Observing from the SGD update~\eqref{eq: SGD} that $\theta_{t+1} - \theta_{t} = -\eta \tgrad L(\theta_t)$, taking the conditional expectation given all randomness up to time $t$, and rearranging terms, 
    % and using the SGD update \eqref{eq: SGD} for $L(\theta)$, 
    we obtain 
    \begin{align*}
        \| \nabla L(\theta_t) \|^2_2  \leq \frac{1}{\eta}(L(\theta_t) -  \E_t [L(\theta_{t+1})])  + \frac{L_N \eta}{2} \E_t \left[\| \tgrad L(\theta_t) \|^2_2  \right].
    \end{align*}
    Taking the full expectation for both sides and using the assumption $\mathbb{E}_t \left[ \left\| \tgrad L(\theta_t) \right\|^2_2 \right] \leq G$ yields
    \begin{align*}
        \E \left[ \| \nabla L(\theta_t) \|^2_2 \right] \leq \frac{1}{\eta}(L(\theta_t) -  \E [L(\theta_{t+1})])  + \frac{L_N  \eta G }{2}.
    \end{align*}
    Finally, we take the average over $ t\in {0, \dots, T - 1}$ and use the telescoping sum to get 
    \begin{align*}
        \frac{1}{T} \sum_{t=0}^{T-1} \mathbb{E} \left[ \left\| \nabla L(\theta_t) \right\|^2_2 \right] \leq \frac{ L(\theta_0) - L^\star}{\eta T} + \frac{L_N  \eta G }{2}.
    \end{align*} 
    Since $\min_{t \in \{0, \dots, T-1\}} \left\| \nabla L(\theta_t) \right\|^2_2 \leq  \frac{1}{T} \sum_{t=0}^{T-1} \left\| \nabla L(\theta_t) \right\|^2_2$, we obtain the first statement of the theorem by replacing the stepsize $\eta = \eta_0 / \sqrt{T}$.
    
    \vspace{12pt}
    \textbf{Gradient bound of infinite-unrolling loss.  }
    For any $\theta$, we consider the norm of the gradient difference:
    \begin{align}
        \|\nabla \cL(\theta) -  \nabla L(\theta) \|_2 &= \left\| \frac{1}{S} \sum_{i=1}^S \nabla_\theta F(\y(\x^\i), \hy(\x^\i)) -  \nabla_\theta F(\y(\x^\i), \hy^K(\x^\i)) \right\|_2 \nonumber\\
        &\leq \frac{1}{S} \sum_{i=1}^S \left\| \nabla_\theta F(\y(\x^\i), \hy(\x^\i)) -  \nabla_\theta F(\y(\x^\i), \hy^K(\x^\i)) \right\|_2, \label{eq: grad difference}
    \end{align}
    which follows from the triangle inequality. We consider one element of the sum, and let $g \coloneqq \left\| \nabla_\theta F(\y(\x), \hy(\x)) -  \nabla_\theta F(\y(\x), \hy^K(\x)) \right\|$. Then, via the chain rule, we have 
    % Note that $\hy(\x) = $
\begin{align*}
    g &= \big\| (J_\theta \y(\x))^\top \nabla_y F(\y(\x), \hy(\x)) + (J_\theta \y(\x))^\top (J_y \FS(\y(\x); \x))^\top \nabla_{\hat{y}} F(\y(\x), \hy(\x)) \\
    &\quad - (J_\theta \y(\x))^\top \nabla_y F(\y(\x), \hy^K(\x)) - (J_\theta \y(\x))^\top (J_y \FS^K(\y(\x);\x))^\top \nabla_{\hat{y}} F(\y(\x), \hy^K(\x)) \big\|_2 \\
    &\leq B_N \big\|\nabla_y F(\y(\x), \hy(\x)) - \nabla_y F(\y(\x), \hy^K(\x)) \big\|_2 \\
    &\quad + B_N \big\| (J_y \FS(\y(\x); \x))^\top \nabla_{\hat{y}} F(\y(\x), \hy(\x)) -  (J_y \FS^K(\y(\x);\x))^\top \nabla_{\hat{y}} F(\y(\x), \hy^K(\x)) \big\|_2.
\end{align*}
where the inequality follows from the bounded Jacobian in Assumption~\ref{assump:nn-bound}.

The first term becomes (by expansion of the gradient terms):
\begin{align*}
    \big\|\nabla_y F(\y(\x), \hy(\x)) - \nabla_y F(\y(\x), \hy^K(\x)) \big\|_2 = \rho \big\| \hy(\x) - \hy^K(\x) \big\|_2.
\end{align*}

The second term:
\begin{align*}
    &\big\| (J_y \FS(\y(\x); \x))^\top \nabla_{\hat{y}} F(\y(\x), \hy(\x)) -  (J_y \FS^K(\y(\x);\x))^\top \nabla_{\hat{y}} F(\y(\x), \hy^K(\x))\big\|_2 \\
    =&\big\| (J_y \FS(\y(\x); \x))^\top \nabla_{\hat{y}} F(\y(\x), \hy(\x)) - (J_y \FS^K(\y(\x);\x))^\top \nabla_{\hat{y}}  F(\y(\x), \hy(\x)) \\
    &+ (J_y \FS^K(\y(\x);\x))^\top \nabla_{\hat{y}}  F(\y(\x), \hy(\x)) -  (J_y \FS^K(\y(\x);\x))^\top \nabla_{\hat{y}} F(\y(\x), \hy^K(\x))\big\|_2
    \\
    \leq & \big\| J_y \FS(\y(\x); \x) - J_y \FS^K(\y(\x);\x)\big\|_2 \big\|\nabla_{\hat{y}} F(\y(\x), \hy(\x))\big\|_2 \\
    &+ \big\| J_y \FS^K(\y(\x);\x)\big\|_2 \big\|\nabla_{\hat{y}}  F(\y(\x), \hy(\x)) -  \nabla_{\hat{y}} F(\y(\x), \hy^K(\x))\big\|_2  \\
    \leq & B_F L_{FS} \big\| \hy(\x) - \hy^K(\x) \big\|_2 +  B_{FS} L_F \big\| \hy(\x) - \hy^K(\x) \big\|_2 \\
    =& (B_F L_{FS} +  B_{FS} L_F) \big\| \hy(\x) - \hy^K(\x) \big\|_2,
\end{align*}
where the second inequality follows from the smoothness of the feasibility-seeking mapping and function $F(\y(\x),\hy(\x))$ (Assumption~\ref{assump:nn-bound}).

Let $\nu = \rho + B_F L_{FS} +  B_{FS} L_F$. Thus, $g$ is upper bounded by
\begin{align*}
    g \leq \nu \| \hy^K(\x) - \hy(\x)\|_2.
\end{align*}

Using this upper bound and Theorem \ref{thm: fs convergence}, we can rewrite the gradient difference \eqref{eq: grad difference} as follows: 
\begin{align}
    \|\nabla \cL(\theta) -  \nabla L(\theta) \|_2   &\leq \frac{1}{S}\sum_{i=1}^S \nu \big\| \hy^K(\x^\i) - \hy(\x^\i) \big\|_2 \nonumber\\
    &\leq \frac{1}{S}\sum_{i=1}^S \frac{8 \nu \eta_\phi L_\phi^2}{\mu_\phi^2} (\phi(\y^\i; \x^\i) - \phi(\hy(\x^\i); \x^\i)), \label{eq: thm1-bound1}
\end{align}
where $\hy(\x^\i) = \FS(\y(\x^\i); \x^\i)$.

Let $ \Phi(\theta) := \frac{1}{S}\sum_{i=1}^S \frac{8 \nu\eta_\phi L_\phi^2}{\mu_\phi^2} \gamma^{K-1} \big(\phi(\y(\x^\i); \x^\i) - \phi(\hy(\x^\i); \x^\i) \big) $.
Using the inverse triangle inequality, the inequality \eqref{eq: thm1-bound1} leads to
\begin{align*}
    \|\nabla \cL(\theta) \|_2 \leq \|  \nabla L(\theta) \|_2 + \Phi(\theta) \gamma^{K-1}.
\end{align*}
We consider the time $t$ and use \eqref{eq: thm2-bound1} to obtain
\begin{align}
    \mathbb{E} [\|\nabla \mathcal{L}(\theta_t) \|_2] &\leq \mathbb{E} [ \|  \nabla L(\theta_t) \|_2] + \mathbb{E} [\Phi(\theta_t)] \gamma^{K-1} \nonumber \\
    &\leq \bigO(T^{-1/4}) +  \mathbb{E} [\Phi(\theta_t)] \gamma^{K-1}. \label{eq: grad bound}
\end{align}
Finally, we need to show that $\mathbb{E} [\Phi(\theta_t)]$ is upper bounded. 
We have 
\begin{align}\label{eq: Phi}
    \mathbb{E} [\Phi(\theta_t)] = \frac{1}{S}\sum_{i=1}^S \frac{8 \nu\eta_\phi L_\phi^2}{\mu_\phi^2} \mathbb{E} \left[\phi \big( \y(\x^\i); \x^\i \big) - \phi \big(\hy(\x^\i); \x^\i \big)\right].
\end{align}
Since Theorem~\ref{thm: fs convergence} shows that the gradient descent converges to the minimum, we have $\phi \big( \y(\x^\i); \x^\i \big) \geq \phi \big(\hy(\x^\i); \x^\i \big)$. Thus, it suffices to show $\E[\phi(\y(\x^\i); \x^\i)]$ is upper bounded for any $i$.
The argument proceeds as follows: We first show that the NN parameters remain bounded during training with SGD, which implies that the network output $\y(\x^\i)$ is also bounded. From there, we conclude that the expectation $\E[\phi(\y(\x^\i); \x^\i)]$ is bounded from above.


\textit{Bound of the NN parameters:} \\
We first establish the bounds of the NN parameters during training with SGD, which via~\eqref{eq: SGD} exhibits
\begin{align*}
    \| \theta_{t+1} - \theta_t \|_2 = \eta \| \tgrad L (\theta_t) \|_2, \quad t =0, \dots, T-1.
\end{align*}
Taking the sum over $j=0, \dots, t-1$ and using triangle inequality and the telescoping sum, we have 
\begin{align*}
     \| \theta_t - \theta_0 \|_2 \leq \sum_{j=0}^{t-1} \left\| \theta_{j+1} - \theta_j\right\|_2 = \sum_{j=0}^{t-1} \eta \| \tgrad L (\theta_j) \|_2.
\end{align*}
Taking the square for both sides and using the Cauchy-Schwarz inequality yields
\begin{align*}
    \| \theta_t - \theta_0 \|^2_2 \leq \sum_{j=0}^{t-1} \eta^2  \sum_{j=0}^{t-1} \| \tgrad L (\theta_j) \|^2_2 \leq \eta^2 T \sum_{j=0}^{t-1} \| \tgrad L (\theta_j)  \|^2_2.
\end{align*}
Using the assumption of bounded variance of the stochastic gradients, we obtain
\begin{align*}
    \E_t [\| \theta_t - \theta_0 \|^2_2] \leq \eta^2 T^2 G.
\end{align*}
By replacing $\eta = \eta_0/\sqrt{T}$ and taking a regular expectation, we finally obtain 
\begin{align}\label{eq: para bound}
    \E [\| \theta_t - \theta_0 \|^2_2] \leq \eta_0^2 T G.
\end{align}

\textit{Bound of the NN output:}

For notational brevity, we denote $y_0$ = $y_{\theta_0} (x)$ and $y_t = y_{\theta_t} (x)$, which are outputs of NNs with initialized parameter $\theta_0$ and parameter $\theta_t$ at time $t$.
By the mean value theorem and the bounded Jacobian $\|J_\theta \y(\x) \|_2 \leq B_N$ in Assumption \ref{assump:nn-bound}, we have 
\begin{align*}
    \|y_t - y_0 \|_2 \leq B_N \| \theta_t - \theta_0 \|_2,
\end{align*}
which, squaring both sides, leads to 
\begin{align}\label{eq: output bound}
    \E[ \|y_t - y_0 \|^2_2] \leq B_N^2 \E [\| \theta_t - \theta_0 \|^2_2] \leq B_N^2 \eta_0^2 TG.
\end{align}

\textit{Bound of $\E[\phi(y_t;\x)]$:}

By the $L$-smoothness of $\phi(y_t; \x)$, we have the quadratic upper bound:
\begin{align*}
    \phi(y_t;\x) \leq \phi(y_0;\x) + \langle \nabla_y \phi(y_0; \x), y_t - y_0 \rangle + \frac{L_\phi}{2} \| y_t - y_0 \|^2_2.
\end{align*}
Taking the expectation on both sides and using \eqref{eq: output bound} yields
\begin{align*}
    \E [\phi(y_t;\x)] \leq \phi(y_0;\x) + B_N \eta \sqrt{TG} \| \nabla_y \phi(y_0; \x)\|_2  + \frac{L_{\phi}}{2} B^2_N \eta_0^2 TG < \infty.
\end{align*}
% which means $\E [\phi(y_t;\x)]$ is upper bounded by a finite value. 
Therefore, we have that $\E[\phi(\y(\x^\i); \x^\i)]$ is upper bounded for any $i$, which implies that $\E[\phi(\hy(\x^\i); \x^\i)]$ is similarly bounded given that $\phi(\hy(\x^\i); \x^\i) \leq \phi(\y(\x^\i); \x^\i)$ by Theorem \ref{thm: fs convergence}
Finally, by \eqref{eq: Phi}, there exists a constant $\kappa >0$ such that $\E[\Phi(\theta_t)] \leq \kappa$ for $t=0, \dots, T-1$. 
Using this in \eqref{eq: grad bound}, we complete the second statement of the theorem. 
\end{proof}

\subsubsection{Proof of Lemma \ref{lemma: stationary prediction}}
\lemmaStationary*
\begin{proof}
    At the stationary parameter $\theta^\star$, the gradient of the loss function vanishes: 
    \begin{align*}
        0 =\nabla \cL(\theta^\star)
        % &= \frac{1}{S} \sum_{i=1}^S \nabla_\theta F(\y(\x^\i), \hy(\x^\i) \\
        = \frac{1}{S} \sum_{i=1}^S (J_\theta y_{\theta^\star}(x^\i))^\top  \bigg(& \nabla_y F(\y(\x^\i), \hy(\x^\i)) \\
        & + (J_y \FS(\y(\x^\i); \x^\i))^\top \nabla_{\hat{y}} F(\y(\x^\i), \hy(\x^\i)) \bigg).
    \end{align*}
    Let $v^\i = \nabla_y F(\y(\x^\i), \hy(\x^\i) + (J_y \FS(y^\i; \x^\i))^\top \nabla_{\hat{y}} F(\y(\x^\i), \hy(\x^\i))$. The stationarity results in
    \begin{align*}
        J(\theta^\star) \left( \begin{array}{c}
              v^{(1)}\\
             \vdots \\
             v^{(S)}
        \end{array} \right) = 0.
    \end{align*}
    Since $J(\theta^\star)$ has full column rank, this implies $(v^{(1)},\dots,v^{(S)})^\top = 0$. 
    Thus, the NN output $y_{\theta^\star}(x^\i)$ and $\hy(\x^\i) = \FS(\y(\x^\i); \x^\i)$ satisfy
    \begin{align*}
        v^\i = \nabla_y F(\y(\x^\i), \hy(\x^\i) + (J_y \FS(y^\i; \x^\i))^\top \nabla_{\hat{y}} F(\y(\x^\i), \hy(\x^\i)) = 0,
    \end{align*}
    which is the stationary condition of $F(\y(\x^\i),\hy(\x^\i))$. 
\end{proof}

\subsubsection{Proof of Theorem \ref{thm: minimizer convergence}}
% \begin{theorem}%\label{thm: minimizer convergence}
%     Given a specific parameter $\x$, suppose that the NN predicts $y_\tau$ which is a global minimizer of $F(y, \hat{y}; \rho_\tau) = f(\hat{y}; x) + \frac{\rho_\tau}{2} \| y - \hat{y}\|^2_2$ with $\hat{y} = \FS(y; x)$, that the mapping $\FS$ is continuous and that $0 < \rho_\tau < \rho_{\tau+1}, \forall \tau$ and $\rho_\tau \rightarrow \infty$. Then, every limit point $(y^\star, \hat{y}^\star)$ of the sequence $\{(y_\tau, \hat{y}_\tau)\}$ admits $y^\star = \hat{y}^\star$, and $\hat{y}^\star$ is a global minimizer of the constrained optimization problem \eqref{prob: nonlinear problem}. 
% \end{theorem}
\thmMinimizer*
\begin{proof}
    Let $\overline{y}$ be a global solution to problem \eqref{prob: nonlinear problem}, that is
    \begin{align*}
        f(\overline{y}; x) \leq f(\hat{y};x)
    \end{align*}
    for all the feasible points $\hat{y}$. We also have that $\overline{y} = \FS(\overline{y}; x)$ given that $\overline{y}$ is already feasible.
    % , $\overline{y}$ is the solution of the feasibility-seeking problem. 
    By the global optimality of $(y_\tau, \hat{y}_\tau)$, we have that $F(y_\tau, \hat{y}_\tau;\rho_\tau) \leq F(\overline{y}, \overline{y};\rho_\tau)$ which results in the inequality 
    \begin{align} \label{eq: proof_ineq_1}
        f(\hat{y}_\tau; x) + \frac{\rho_\tau}{2} \|y_\tau - \hat{y}_\tau \|^2_2 \leq f(\overline{y}; x) + \frac{\rho_\tau}{2} \|\overline{y} - \overline{y}\|^2_2 = f(\overline{y}; x).
    \end{align}
    By rearranging this inequality, we obtain 
    \begin{align} \label{eq: proof_ineq_2}
        \| y_\tau - \hat{y}_\tau\|^2_2 \leq \frac{2}{\rho_\tau}(f(\overline{y}; x) - f(\hat{y}_\tau; \x)).
    \end{align}
    Suppose that $y^\star$ is a limit point of $\{y_\tau\}_{\tau\geq 0}$, so that there is an infinite subsequence $\mathcal{T}$ such that 
    \[
    \lim_{\tau \in \mathcal{T}} y_\tau = y^\star.
    \]
    As the feasibility seeking mapping is continuous by assumption, we can define $\hat{y}^\star \coloneqq \lim_{\tau \in \cT} \mathsf{FS}(y_\tau;\x)$. By the continuity, we have 
    \[
    \lim_{\tau \in \mathcal{T}} \| y_\tau - \hat{y}_\tau\|^2_2 = \lim_{\tau \in \mathcal{T}} \| y_\tau - \FS(y_\tau; x)\|^2_2 = \| y^\star - \hat{y}^\star \|^2_2.
    \]  
    Taking the limit as $\tau \rightarrow \infty, \tau \in \mathcal{T}$ on both sides of \eqref{eq: proof_ineq_2} yields
    \begin{align*}
        \| y^\star - \hat{y}^\star\|^2_2 = \lim_{\tau \in \mathcal{T}} \| y_\tau - \hat{y}_\tau\|^2_2 \leq \lim_{\tau \in \mathcal{T}} \frac{2}{\rho_\tau}(f(\overline{y}; x) - f(\hat{y}_\tau; \x)) = 0.
    \end{align*}
    %By the continuity of the norm
    Therefore, we have that $y^\star = \hat{y}^\star$. Note that $\hat{y}^\star$ is the output of the feasibility-seeking mapping, so $y^\star$ and $ \hat{y}^\star$ are feasible. 
    Moreover, by taking the limit as $\tau \rightarrow \infty, \tau \in \mathcal{T}$ in \eqref{eq: proof_ineq_1}, we have 
    \begin{align*}
        f(\hat{y}^\star) \leq f(\hat{y}^\star) + \lim_{\tau \in \mathcal{T}} \frac{\rho_\tau}{2} \| y_\tau - \hat{y}_\tau \|^2_2 \leq f(\overline{y}; x).
    \end{align*}
    Since $\hat{y}^\star$ is a feasible point with an objective value no larger than that of the global solution $\overline{y}$, it must itself be a global solution.
\end{proof}
Using a similar argument, we can show that \( \hat{y}^\star \) is a minimizer of Problem~\ref{prob: nonlinear problem} under the same assumptions, even when \( \rho = 0 \). However, in practice, we often choose \( \rho > 0 \) for the reasons discussed in Section~\ref{sec: fs} and Appendix~\ref{appendix: more experiments}.


\subsection{Analysis of Truncated Unrolled Differentiation}\label{subsec: truncated diff analysis}
In this section, we build intuition for why truncating backpropagation remains effective in practice. Specifically, if the violation function \(\phi(s,x)\) is locally strongly convex, the bias introduced by truncation decays exponentially with the truncation depth \(K'\). Consequently, a relatively small \(K'\) suffices to yield a high-quality gradient estimate.


\begin{assumption}\label{assum: local strongly convex}
    The violation function $s \mapsto \phi(s;\x)$ is twice differentiable. For $s_K \in \R^n$, there exists a neighborhood $\cN$ of $s_K$ where $s \mapsto \phi(s; \x)$ is $\mu'$-strongly convex. 
    Moreover, for some $K_0 \in \N$, the iterates satisfy $s_k \in \cN$ for $k \geq K_0$.
\end{assumption} 
Although the violation function involves a $\operatorname{max}$ operation and is not twice differentiable at the kink points of the $\operatorname{max}$, we can approximate the $\operatorname{max}$ in practice using the softplus function, which is infinitely differentiable and can approximate the $\operatorname{max}$ arbitrarily well. 
Therefore, this assumption remains reasonable in practice due to the smooth approximation provided by the softplus function.
This assumption enables us to bound the bias caused by the truncation as in the following lemma. 

\begin{lemma}\label{lemma: Jac error}
    Suppose Assumption \ref{assum: local strongly convex} holds and that the feasibility-seeking procedure is unrolled with the stepsize $\eta_\phi \in (0, 1/L_\phi]$ for a total of $K$ iterations, but gradient tracking is performed only for the first $K' \in [K_0, K]$ iterations, with the remaining iterations treated as pass-through during backpropagation. Define the true Jacobian when all $K$ iterations are tracked as $J_y \FS^K(\y(\x); \x)$ and the truncated Jacobian as $J_y^\mathrm{trun} \FS^K(\y(\x);\x)$: 
    \begin{align*}
        J_y \FS^K(\y(\x); \x)  = \frac{\partial s_K}{\partial s_0}, \qquad
        J_y^\mathrm{trun} \FS^K(\y(\x);\x)  = \frac{\partial s_{K'}}{\partial s_0}        
    \end{align*}    
    Then, the bias caused by the truncation satisfies 
    \begin{align*}
        \left\| J_y \FS^K(\y(\x); \x) -  J_y^\mathrm{trun} \FS^K(\y(\x);\x) \right\|_2 \leq C(1 - \delta^{K-K'}) \delta^{K'}
    \end{align*}
    where $\delta = 1- \eta_\phi \mu' \in [0,1)$ and $C$ is a finite positive constant.
\end{lemma}

Since the truncation occurs during the backward pass, it also perturbs the stochastic gradient. 
We denote the resulting truncated gradient estimate by $\tgrad^\mathrm{trun} L(\theta_t)$.
The next lemma quantifies this effect.
\begin{lemma} \label{lemma: bounded grad est error}
    For gradient estimate without truncation, suppose $\mathbb{E}_t \left[ \| \tgrad L(\theta_t) \|^2_2 \right] \leq G$ and $\E_t[\tgrad L(\theta_t)] = \nabla L(\theta_t)$ for all $t$. Then, when the gradient truncation is applied, under Assumptions in Lemma \ref{lemma: Jac error}, the truncated gradient estimate satisfies 
    \begin{align*}
        &\E_t \left[\| \tgrad^\mathrm{trun} L(\theta_t) \|^2_2\right] \leq G + \sigma_{K'}, \\
        &\E_t\left[ \tgrad^\mathrm{trun} L(\theta_t) \right] = \nabla L(\theta_t) + \varepsilon_{K'}.
    \end{align*}
    where $\sigma_{K'} = \bigO(\delta^{K'})$ and $\varepsilon_{K'} \in \R^{n_\theta}$ with $\|\varepsilon_{K'} \|= \bigO(\delta^{K'})$.
\end{lemma}

This lemma shows that the truncation introduces bounded biases to the variance and expectation of the gradient estimate, yet these biases decay exponentially fast in the truncation depth \(K'\).
This qualification allows us to understand the behavior of the SGD when the truncation is applied.
\begin{theorem}\label{thm: SGD convergence 2}
    Suppose Assumption~\ref{thm: fs convergence}, \ref{assump:nn-bound}, \ref{assump:nn-smooth}, and the assumptions of Lemma~\ref{lemma: bounded grad est error} hold.
    Then, after running the SGD for $T$ iterations with step size $\eta = \eta_0 / \sqrt{T}$ with $\eta_0 > 0$, there is at least one iteration $t \in \{0, \dots, T-1\}$ satisfying 
    \begin{align}%\label{eq: thm2-bound1}
        &\mathbb{E} \left[ \left\| \nabla L(\theta_t) \right\|^2_2 \right] \leq \bigO\left( T^{-1/2} + \delta^{K'} \right), \\
        &\mathbb{E} [\|\nabla \cL(\theta_t) \|_2] \leq 
        \bigO\left( T^{-1/4} + \gamma^K + \delta^{K'} \right),
    \end{align}
    where $\gamma$ and $\delta$ are defined as in Theorem \ref{thm: SGD convergence 2} and Lemma \ref{lemma: Jac error}.
\end{theorem}

This theorem characterizes how the bias caused by truncating the gradient tracking in the feasibility-seeking procedure propagates to the convergence of SGD. 
We see that the upper bounds of the expected gradient norms contain bias terms of order $\bigO(\gamma^K)$ and $\bigO(\delta^{K'})$. 
Since $0 < \gamma, \delta \leq 1$, even modest values of $K$ and $K'$ suffice to render these biases negligible in practice.
We now present the proofs of the above lemmas and theorems.

\subsubsection{Proof of Lemma \ref{lemma: Jac error}}
Lemma \ref{lemma: Jac error} is an immediate result of the following lemma.
\begin{lemma}\label{lemma: general Jac error}
    Let $\varphi: \R^d \rightarrow \R$ be twice differentiable, $L$-smooth, satisfy the PL condition with constant $\mu$. Suppose there exists a neighborhood $\cN$ of $s_K$ on which $\varphi$ is strongly convex, i.e., $\nabla^2 \varphi(s) \succeq \mu' I, \forall s \in \cN$ with $\mu' > 0$. Define 
    \begin{align*}
        J^\mathrm{true} = \frac{\partial s_K}{ \partial s_0} = \prod_{k=0}^{K-1} (I - \eta \nabla^2 \varphi (s_k)),
        \quad 
        J^\mathrm{trun} = \frac{\partial s_{K'}}{ \partial s_0} = \prod_{k=0}^{K'-1} (I - \eta \nabla^2 \varphi (s_k)).
    \end{align*}
    If $s_k \in \cN$ for all $k \geq K_0$ and stepsize of gradient descent is chosen $\eta \in (0, {1}/{L}]$, then for any $K' \in [K_0, K]$, the error matrix $E = J^\mathrm{true} - J^\mathrm{trun}$ satisfies
    \begin{align*}
        \| E \|_2 \leq C ( 1 - \delta^{K-K'}) \delta^{K'},
    \end{align*}
    where $\delta = 1 - \eta \mu' \in [0, 1)$ and $C = \frac{L}{\mu'} \delta^{-K_0}(1 + \eta L)^{K_0}$.
\end{lemma}
\begin{proof}
    Define $J_k = \frac{\partial s_{k+1}}{\partial s_k} = I - \eta \nabla^2 \varphi(s_k)$ and rewrite the norm of error matrix $E$ as 
    \begin{align}\label{eq: trunc E}
        \| E\|_2 &= \|(J_{K-1} \dots J_{K'} - I)(J_{K'-1} \dots J_0) \|_2 \nonumber\\
        &\leq \| J_{K-1} \dots J_{K'} - I \|_2 \| J_{K'-1} \dots J_0\|_2.
    \end{align}
    %%% 
    \textbf{Step 1:} $0 \leq k < K'$. \\
    For $k < K_0$, the $L$-smoothness gives $-L I \preceq \nabla^2 \varphi(s_k) \preceq LI$, so 
    \begin{align*}
        \| J_k \|_2 = \| I - \eta \nabla^2 \varphi(s_k) \|_2 \leq 1 + \eta L.
    \end{align*}
    For $K_0 \leq k <K'$, strong convexity implies $\mu' I \preceq \nabla^2 \varphi (s_k) \preceq L I$, and then we obtain 
    \begin{align*}
       \| J_k \|_2 \leq 1 - \eta \mu'.
    \end{align*}
    Let $\delta := 1 - \eta \mu' \in [0, 1)$.
    Then, we use the sub-multiplicativity property of the spectral norm to obtain
    \begin{align}\label{eq: trunc B1}
        \| (J_{K'-1} \dots J_0)\|_2 \leq \| (J_{K'-1} \dots J_{K_0})\|_2 \| (J_{K_0-1} \dots J_0)\|_2 \leq \delta^{K' - K_0} (1 +\eta L)^{K_0}.
    \end{align}
    
    \textbf{Step 2:} $K' \leq k < K$. \\
    Use the telescoping sum for $\| (J_{K-1} \dots J_{K'} - I) \|_2$: 
    \begin{align*}
        \| J_{K-1} \dots J_{K'} - I \|_2 &=  \left \|\sum_{k = K'}^{K-1} \left( \prod_{j=k+1}^{K-1} J_j \right) (J_k-I) \right\|_2 \\
        & \leq \sum_{k = K'}^{K-1} \left\|\prod_{j=k+1}^{K-1} J_j \right\|_2 \left\| J_k-I \right\|_2.
    \end{align*}
    From the $L$-smoothness we have $\| J_k - I \|_2 = \| \eta \nabla^2 \varphi (s_k) \|_2 \leq \eta L$. 
    For $k \geq K' > K_0$, by assumption, we are in the neighborhood $\cN$ where $\varphi$ is locally strongly convex, and thus we obtain 
    \begin{align*}
        \left\Vert \prod_{j=k+1}^{K-1} J_j \right\Vert_2 \leq \delta^{K-k-1},
    \end{align*}
    which results in 
    \begin{align*}
        \| J_{K-1} \dots J_{K'} - I \|_2 &\leq \eta L \sum_{k = K'}^{K-1} \delta^{K-k-1}.      
    \end{align*}
    We change the index of the summation to $i=K-k-1$. When $k=K'$, $i=K-K'-1$. When $k=K-1$, $i=0$. The sum becomes a finite geometric series with $0 \leq \delta < 1$:
    \begin{align}\label{eq: trunc B2}
        \| J_{K-1} \dots J_{K'} - I \|_2 &\leq \eta L \sum_{i = 0}^{K-K'-1} \delta^{i} \nonumber\\
        &= \eta L\frac{1 - \delta^{K-K'}}{1- \delta}   \nonumber\\
        &= \frac{L} {\mu'}(1 - \delta^{K-K'}).
    \end{align}

    \textbf{Step 3:} Combining the bounds  \\
    Substituting \eqref{eq: trunc B1} and \eqref{eq: trunc B2} into the inequality of error matrix \eqref{eq: trunc E} yields
    \begin{align*}
        \| E\|_2 \leq \frac{L}{\mu'}(1- \delta^{K-K'}) \delta^{K' - K_0} (1 +\eta L)^{K_0}.
    \end{align*}
    Define the constant $C = \frac{L}{\mu'} \delta^{-K_0}(1 + \eta L)^{K_0}$. We can rewrite the bound as follows:
    \begin{align*}
        \| E\|_2 \leq C (1- \delta^{K-K'})\delta^{K'}.
    \end{align*}    
\end{proof}


\subsubsection{Proof of Lemma \ref{lemma: bounded grad est error}}

\begin{proof}
    For a minibatch with $D$ samples, we have
    \begin{align}
        &\left\| \tgrad L(\theta_t) - \tgrad^\mathrm{trun} L(\theta_t) \right\|_2 \nonumber\\
        &= \Bigg\| \frac{1}{D} \sum_{i=1}^D
         (J_\theta \y (x^\i))^\top \nabla_y F(\y(\x^\i), \hy^K(\x^\i)) \nonumber\\
         &\qquad + (J_\theta \y (x^\i))^\top (J_y \FS^K(\y(\x)^\i; \x^\i))^\top \nabla_{\hat{y}} F(\y(\x^\i), \hy^K(\x^\i))  \nonumber\\
         &\qquad - (J_\theta \y (x^\i))^\top \nabla_y F(\y(\x^\i), \hy^K(\x^\i)) \nonumber\\
         &\qquad - (J_\theta \y (x^\i))^\top (J_y^\trun \FS^K(\y(\x)^\i; \x^\i))^\top \nabla_{\hat{y}} F(\y(\x^\i), \hy^K(\x^\i)) \Bigg\|_2 \nonumber\\
         &= \Bigg\| \frac{1}{D} \sum_{i=1}^D  (J_\theta \y (x^\i))^\top (J_y \FS^K(\y(\x)^\i; \x^\i) \nonumber\\
         &\qquad - J_y^\trun \FS^K(\y(\x)^\i; \x^\i))^\top \nabla_{\hat{y}} F(\y(\x^\i), \hy^K(\x^\i)) \Bigg\|_2 \nonumber\\
         &\leq \frac{1}{D} \sum_{i=1}^D B_N B_F C (1 - \delta^{K-K'})\delta^{K'} \qquad \text{(by Lemma \ref{lemma: Jac error})} \nonumber \\
         &= B_N B_F C (1 - \delta^{K-K'})\delta^{K'}. \label{eq: estimated grad bound}
    \end{align}
    Then, by triangle inequality, we have the following bound 
    \begin{align*}
        \E_t \left[ \left\| \tgrad^\mathrm{trun} L(\theta_t) \right\|^2_2\right] &\leq \E_t \left[ \left(\left\|  \tgrad^\mathrm{trun} L(\theta_t) - \tgrad L(\theta_t)  \right\|_2 + \left\| \tgrad L(\theta_t) \right\|_2 \right)^2 \right] \\
        &\leq B_N B_F C (1 - \delta^{K-K'})^2\delta^{2K'} + G + 2 B_N B_F C\sqrt{G}(1 - \delta^{K-K'})\delta^{K'} \\
        &= G + \bigO(\delta^{K'}).
    \end{align*}

    In addition, by \eqref{eq: estimated grad bound}, we have that 
    \begin{align*}
       \left\| \tgrad L(\theta_t) - \tgrad^\mathrm{trun} L(\theta_t) \right\|_\infty \leq  \left\| \tgrad L(\theta_t) - \tgrad^\mathrm{trun} L(\theta_t) \right\|_2 \leq B_N B_F C (1 - \delta^{K-K'})\delta^{K'} \coloneqq \upsilon_{K'},
    \end{align*}
    which results in 
    \begin{align*}
        \tgrad L(\theta_t) - \upsilon_{K'} \mathbf{1} \leq \tgrad^\mathrm{trun} L(\theta_t) \leq \tgrad L(\theta_t) + \upsilon_{K'} \mathbf{1} \quad \text{(element-wise)}.
    \end{align*}
    There exists $\varepsilon_{K'} \in \R^{n_\theta}$ such that $\|\varepsilon_{K'}\|_2 \leq \upsilon_{K'} \sqrt{n_\theta}$ that satisfies 
    \begin{align*}
        \E_t\left[ \tgrad^\mathrm{trun} L(\theta_t) \right] = \E_t\left[ \tgrad L(\theta_t) \right] + \varepsilon_{K'} = \nabla L(\theta_t) + \varepsilon_{K'}.
    \end{align*}    
    Since $v_{K'} = \bigO(\delta^{K'})$, the proof is complete.
\end{proof}

\subsubsection{Proof of Theorem \ref{thm: SGD convergence 2}}
\begin{proof}
    \textbf{Gradient bound for finite-unrolling loss with truncated gradient.}
    By the quadratic upper bound of the $L$-smooth function and bounded variance of the stochastic gradient (Lemma~\ref{lemma: bounded grad est error}):
    \begin{align}\label{eq: trunc_thm 1}
        \E_t [L(\theta_{t+1})] &\leq L(\theta_t) -\eta \langle \nabla L(\theta_t), \E_t[\tgrad^\trun L(\theta_t)] \rangle + \frac{L_N \eta^2}{2} \E_t \left[\| \tgrad^\trun L(\theta_t) \|^2_2\right] \nonumber\\
        &= L(\theta_t) -  \eta \langle \nabla L(\theta_t), \nabla L(\theta_t) + \varepsilon_{K'} \rangle + \frac{L_N \eta^2}{2} (G +\sigma_{K'}).
    \end{align}
    We observe that 
    \begin{align*}
        \langle \nabla L(\theta_t), \nabla L(\theta_t) + \varepsilon_{K'} \rangle
        &= \|\nabla L(\theta_t) \|^2_2 + \langle  \nabla L(\theta_t), \varepsilon_{K'} \rangle + \frac{1}{2}\|\varepsilon_{K'} \|^2_2 - \frac{1}{2}\|\varepsilon_{K'}\|^2_2\\
        &= \frac{1}{2}\|\nabla L(\theta_t) \|^2_2 - \frac{1}{2}\|\varepsilon_{K'}\|^2_2 + \frac{1}{2} \left(\|\nabla L(\theta_t) \|^2_2 + 2\langle  \nabla L(\theta_t), \varepsilon_{K'} \rangle +\|\varepsilon_{K'}\|^2_2 \right) \\
        &= \frac{1}{2}\|\nabla L(\theta_t) \|^2_2 - \frac{1}{2}\|\varepsilon_{K'}\|^2_2 + \frac{1}{2}\left\| \nabla L(\theta_t) + \varepsilon_{K'} \right\|^2_2\\
        &\geq \frac{1}{2}\|\nabla L(\theta_t) \|^2_2 - \frac{1}{2}\|\varepsilon_{K'}\|^2_2.
    \end{align*}
    Using this in $\eqref{eq: trunc_thm 1}$ and rearranging terms, we obtain 
    \[
    \| \nabla L(\theta_t) \|^2_2 \leq \frac{2}{\eta} (L (\theta_t) - \E_t[L(\theta_{t+1}]) + \| \varepsilon_{K'} \|^2_2 + L_N \eta (G + \sigma_{K'}).
    \]
    Taking the expectation for both sides and the average over $T$ steps, we obtain
    \begin{align*}
        \frac{1}{T}\sum_{t=0}^{T-1} \E [\| \nabla L(\theta_t) \|^2_2] &\leq \frac{2}{\eta T} \E \left[ \sum_{t=0}^{T-1} L(\theta_t) - L(\theta_{t+1}) \right] + \| \varepsilon_{K'} \|^2_2 + L \eta (G + \sigma_{K'}) \\
        &\leq \frac{2}{\eta T} \left(L(\theta_0) - L^\star \right) + \| \varepsilon_{K'} \|^2_2 + L_N \eta (G + \sigma_{K'}).
    \end{align*}
    Since $\min_{t \in \{0, \dots, T-1\}} \left\| \nabla L(\theta_t) \right\|^2_2 \leq  \frac{1}{T} \sum_{t=0}^{T-1} \left\| \nabla L(\theta_t) \right\|^2_2$, we obtain the first statement of the theorem by replacing the stepsize $\eta = \eta_0 / \sqrt{T}$.

    Using the same argument as the proof of Theorem \ref{thm: SGD convergence 1}, there exists positive constant $\kappa_1, \kappa_2$ such that
    \begin{align*}
        \mathbb{E} [\|\nabla \cL(\theta_t) \|_2] \leq \kappa_1 T^{-1/4} +  \kappa_2 \gamma^{K-1} + \| \varepsilon_{K'} \|^2_2 + \frac{L_N \eta_0}{\sqrt{T}} \sigma_{K'}.
    \end{align*}
\end{proof}







\end{document}
